\newfontfamily\BBFifthTranslit[Path=vendor/fonts/,BoldFont=DejaVuSerif-Bold.ttf,ItalicFont=DejaVuSerif.ttf,BoldItalicFont=DejaVuSerif-Bold.ttf]{DejaVuSerif.ttf}
\chapter{LINGUISTIQUE COMPARATIVE ET LINGUISTIQUE TYPOLOGIQUE}
\label{chapter-5}
\markright{Linguistique comparative et typologique}
\origpage{131}{117}
\begin{keybox}[Plan du chapitre]
\small
\textbf{I.\quad LA LINGUISTIQUE COMPARATIVE}
\begin{enumerate}[label=\arabic*.]
\item Les pionniers du comparatisme
\begin{enumerate}[label=\textcircled{\scriptsize\arabic*}]
\item Marcus van Boxhorn
\item Gaston-Laurent Cœurdoux
\item William Jones
\item La notion d’hyperpolyglossie
\end{enumerate}
\item Les fondateurs d’une nouvelle science
\begin{enumerate}[label=\textcircled{\scriptsize\arabic*}]
\item Frédéric Schlegel
\begin{enumerate}[label=\arabic*.]
\item \textit{Sur la langue et la sagesse des Indiens} (1808)
\item Comparatisme et anatomie
\end{enumerate}
\item Rasmus Rask
\item Franz Bopp
\item August Schleicher
\end{enumerate}
\end{enumerate}
\textbf{II.\quad LE CLASSEMENT GÉNÉTIQUE DES LANGUES}
\begin{enumerate}[label=\arabic*.]
\item Sept mille langues : comment les classer ?
\begin{enumerate}[label=\textcircled{\scriptsize\arabic*}]
\item La notion de famille de langues
\item La notion de branche
\item La notion de superfamille
\item La notion d’isolat linguistique
\item Les langues sinoxéniques
\end{enumerate}
\item La famille indo-européenne
\begin{enumerate}[label=\textcircled{\scriptsize\arabic*}]
\item La branche romane
\begin{enumerate}[label=\arabic*.]
\item Diglossie latine
\begin{enumerate}[label=\arabic*)]
\item Du \textit{sermo urbanus} au \textit{sermo plebeius}
\item Valérius Probus : l’« Appendix Probi »
\end{enumerate}
\item Naissance de la langue romane
\begin{enumerate}[label=\arabic*)]
\item Le concile de Tours (813)
\item Les \textit{Serments de Strasbourg} (842)
\end{enumerate}
\item Variété des langues romanes
\origpage{132}{118}
\begin{enumerate}[label=\arabic*)]
\item La notion de substrat
\item La notion de superstrat
\end{enumerate}
\item Les langues d’oc et langues d’oïl
\end{enumerate}
\item La branche germanique
\item La branche celtique
\item La branche indo-iranienne
\item Les branches éteintes
\begin{enumerate}[label=\arabic*.]
\item La branche anatolienne
\item La branche tokharienne : koutchéen et agnéen
\end{enumerate}
\end{enumerate}
\end{enumerate}
\textbf{III.\quad LA LINGUISTIQUE TYPOLOGIQUE}
\begin{enumerate}[label=\arabic*.]
\item Les grands contributeurs
\begin{enumerate}[label=\textcircled{\scriptsize\arabic*}]
\item Auguste Schlegel : la tripartition des langues
\item Wilhelm von Humboldt
\begin{enumerate}[label=\arabic*.]
\item L’hypothèse humboldtienne
\item Le chinois, « une langue sans grammaire » ?
\begin{enumerate}[label=\arabic*)]
\item La réponse d’Abel Rémusat
\item Le « génie de la langue chinoise »
\end{enumerate}
\end{enumerate}
\item Joseph Harold Greenberg
\begin{enumerate}[label=\arabic*.]
\item Les invariants linguistiques
\item Typologie syntaxique
\begin{enumerate}[label=\arabic*)]
\item Les langues SOV, SVO et VSO
\item Les langues VOS, OVS et OSV
\end{enumerate}
\end{enumerate}
\end{enumerate}
\item Typologie morphologique
\begin{enumerate}[label=\textcircled{\scriptsize\arabic*}]
\item Les langues isolantes
\item Les langues flexionnelles
\begin{enumerate}[label=\arabic*.]
\item Les langues agglutinantes
\begin{enumerate}[label=\arabic*)]
\item Le \textit{kiswahili}
\item Le finnois
\item Le tamoul
\item Le turc
\end{enumerate}
\item Les langues synthétiques
\begin{enumerate}[label=\arabic*)]
\item La notion de paradigme
\item Différences entre langues agglutinantes et synthétiques
\item Illustration : illatif et inessif
\end{enumerate}
\end{enumerate}
\end{enumerate}
\end{enumerate}
\textbf{Exercices}\par
\textbf{Pour aller plus loin}
\end{keybox}

\origpage{133}{119}
\section*{Introduction}
Sous l’effet d’une double influence, le début du 19\textsuperscript{e} siècle a été témoin d’une importante mutation dans la réflexion sur le langage.

D’une part, dans un contexte de réveil des nationalismes, de prise de conscience des passés nationaux et d’exaltation de l’âme profonde des peuples, le langage a été considéré comme une antiquité nationale et un produit de l’histoire. Cette période a donc connu une intense activité philologique portant sur la littérature populaire : contes, légendes, épopées, mythologie, etc.

D’autre part, parallèlement au romantisme, les sciences de la nature ont exercé une influence opposée mais tout aussi capitale sur l’étude du langage. En plein essor à cette époque, les idées du naturaliste français Jean-Baptiste Lamarck (1744-1829), du zoologue et paléontologue Georges Cuvier (1769-1832), ainsi que la théorie de l’évolutionnisme de Charles Darwin (1809-1882) ont commencé à pénétrer le champ des études sur le langage.\ednote{年代范围须区别：段首说19世纪初，后文原书第125页自身却把Schleicher与Darwin理论的联系放在1860年代。本段把不同时段的影响压缩在同一时期，宜把有关影响的年代分别交代。}

Sous cette double influence de l’histoire et du naturalisme, les linguistes se sont lancés à la recherche de la « langue mère » et ont cherché à regrouper les langues du monde en « familles » partageant une descendance commune. En faisant cela, ils ont inventé une nouvelle science : la grammaire comparée, considérée comme la première école de linguistique moderne dotée d’une méthode scientifique.

\section{LA LINGUISTIQUE COMPARATIVE}
Si la linguistique comparée s’est développée à partir du 19\textsuperscript{e} siècle, la ressemblance entre les langues avait été remarquée par des observateurs de période bien plus ancienne comme le poète italien \textbf{Dante} (1265-1321) qui, dans \textit{De l’éloquence en langue vulgaire} (1304), faisait état de ressemblances entre l’italien, l’espagnol et le provençal. Mais il faudra encore attendre six cents ans\ednote{原文如此。按本页所列1304与1816计算，相隔512年，应为约五百年而非六百年。} et la découverte du sanskrit, la langue ancienne de l’Inde, pour assister à la naissance d’une véritable science. Pour \textbf{Franz Bopp} (1791-1867), le fondateur du comparatisme, la découverte du sanskrit a été « comme la découverte d’un nouveau monde ».

\subsection{Les pionniers du comparatisme}
Si l’allemand Franz Bopp est généralement reconnu pour avoir fondé en 1816 la grammaire comparée, la contribution de ses prédécesseurs un peu oubliés avec le temps est loin d’être négligeable.

\minititle{\textcircled{\scriptsize 1} Marcus van Boxhorn}
\textbf{Marcus Zuerius van Boxhorn} (1612-1653), homme politique et linguiste néerlandais, est le premier à avoir soupçonné l’existence d’une ancienne langue commune aux langues néerlandaise,
\origcont{134}{120}grecque, latine, perse, germaniques, slaves, celtes et baltes. Il avait baptisé cette langue-ancêtre du nom « scythique », d’après le nom des Scythes, peuple nomade des steppes eurasiatiques. Boxhorn avait publié une première lettre savante intitulée « Des mots perses enregistrées\ednote{原文如此。分词修饰阳性复数mots，应作enregistrés。} par Quinte Curce et de leur parenté avec des termes germaniques ». Mais c’est la découverte en 1645, dans la province de Zélande aux Pays-Bas, de stèles romaines provenant du sanctuaire de Nehalennia, une déesse celto-germanique, qui semble avoir provoqué l’intuition de Boxhorn. Cherchant à ramener le nom de la déesse à une racine commune à plusieurs langues européennes, il a émis l’hypothèse d’une origine commune et l’a exposée dans une lettre savante intitulée « Éclaircissements sur la déesse Nehalennia ». Une telle recherche s’inscrivait dans la quête d’une langue originelle, d’une « langue mère » pour l’humanité. À cette époque, la genèse des langues à partir de l’hébreu comme langue mère, idée dominante jusqu’à la Renaissance, commençait à être contestée par d’autres modèles remettant en cause l’idée d’une origine unique. Boxhorn appuyait son intuition sur une démonstration systématique lexicale et morphologique, mais ses travaux n’ont pas joui d’une grande postérité : ses écrits, publiés en néerlandais, ont eu une diffusion limitée en Europe.

\minititle{\textcircled{\scriptsize 2} Gaston-Laurent Cœurdoux}
Missionnaire jésuite français, \textbf{Gaston-Laurent Cœurdoux}(1691-1779) est surtout connu aujourd’hui comme spécialiste de la culture et des langues de l’Inde. Son \textit{Dictionnaire trilingue télougou (une des 234 langues de l’Inde) français et sanscrit} fait encore autorité aujourd’hui. Dans un mémoire envoyé en 1767 à l’Académie des inscriptions et belles-lettres, Cœurdoux a démontré l’analogie existant entre le sanscrit, le latin, le grec, l’allemand et le russe. Au lendemain de la Révolution française, indianiste\ednote{原文如此。此处名词短语缺限定词，应为l’indianiste français。} français \textbf{Abraham Hyacinthe Anquetil-Duperron} a fait publier un chapitre du mémoire de Cœurdoux. En 1808, le missionnaire indianiste \textbf{Jean-Antoine Dubois} (1766-1848) a lui utilisé les travaux de Cœurdoux... mais en les faisant passer comme étant son propre ouvrage. Ce n’est que récemment que la vérité historique a été rétablie, et l’on reconnaît à Cœurdoux la paternité de la découverte du lien entre le sanscrit et les langues anciennes européennes.

\minititle{\textcircled{\scriptsize 3} William Jones}
Né à Londres et diplômé d’Oxford, l’orientaliste \textbf{William Jones} (1746-1794), publie en 1771 une \textit{Dissertation sur la littérature orientale} et une \textit{Grammaire de la langue persane}. Ces activités n’étant pas suffisamment rémunératrices, il se consacre, à partir de 1770 et durant trois années, à l’étude du droit. En 1783, il est nommé juge à la cour suprême de Calcutta, ce qui entraîne son départ et son installation dans les Indes orientales. À son arrivée, il tombe sous le charme de la culture indienne et fonde la Société asiatique du Bengale. Au moment de sa mort à Calcutta en 1794, William Jones maîtrisait treize langues (l’hébreu, le français, le chinois, le sanskrit, etc.) et en connaissait vingt-huit autres. Les publications de la Société asiatique du Bengale ont rendu le sanskrit accessible aux chercheurs européens et ont grandement contribué au développement de la philologie comparative. Dans son \textit{Troisième discours} à la Société asiatique de Calcutta, en 1786, Jones a écrit :

\origpage{135}{121}
\begin{lecture}{}
La langue sanskrite, quelle que soit son antiquité, est d’une structure admirable ; plus parfaite que la grecque ; plus ample que la latine, et plus exquisément raffinée qu’aucune des deux mais ayant envers chacune d’entre elles deux une affinité plus forte, tant dans les racines des verbes que dans les formes de la grammaire, qu’il n’en pourrait avoir résulté par accident ; si forte en vérité qu’aucun philologue ne les pourrait examiner toutes trois sans croire qu’elles ont surgi de quelque source commune, qui, peut-être, n’existe plus. Il y a du reste une raison similaire, quoique pas tout à fait aussi contraignante, pour supposer que le gotique et le celtique, s’ils ont été mêlés par la suite avec un parler différent, n’en descendent pas moins en définitive de la même origine que le sanskrit ; on pourrait ajouter en outre à cette famille le vieux perse.
\end{lecture}

\minititle{\textcircled{\scriptsize 4} La notion d’hyperpolyglossie}
Tous les comparatistes que nous venons de présenter, et ceux qu’on va voir plus bas, ont la particularité enviable de connaître des dizaines de langues : ce sont des hyperpolyglottes, c’est-à-dire des personnes qui parlent couramment six langues ou davantage. Ce terme, inventé en 2003 par le linguiste anglais \textbf{Richard Anthony Hudson} (1939-) est à rapprocher du terme polyglotte, désignant une personne parlant plusieurs langues.

Plusieurs théories ont été développées pour expliquer pourquoi certaines personnes peuvent aisément apprendre de nombreuses langues, alors que la plupart des gens ne peuvent en apprendre que quelques-unes, et avec tellement de difficultés. En 2004, des scientifiques ont publié une analyse du cerveau de l’hyperpolyglotte allemand \textbf{Emil Kreb} (1867-1930)\ednote{原文此处作Kreb；下文及所引论文均作Krebs，应补末尾s。} qui maîtrisait soixante-huit langues à l’écrit comme à l’oral et en avait étudié également cent-vingt autres. Peu après la mort de Krebs, son cerveau a été transmis à l’Institut de recherche sur le cerveau Kaiser Wilhelm à Berlin. En 2004, trois scientifiques, Katrin Amunts, Axel Schleicher et Karl Zilles ont publié \textit{Outstanding language competence and cytoarchitecture in Broca’s speech region}. Leur analyse de la composition cellulaire (cytoarchitectonie) du cerveau de Krebs, en particulier la répartition des différents types de cellules, leur nombre et leur distribution spatiale, a révélé des changements dans la construction de son aire de Broca, qui est, on le sait, responsable du traitement du langage :

\begin{lecture}{}
We analyze the cytoarchitecture of areas 44 and 45 (anatomical correlates of Broca’s speech region) of a person with a documented extraordinary competence in language performance (Emil Krebs, E.K.), and compared it with 11 control brains. Morphometry and multivariate statistical analysis revealed significant cytoarchitectonic differences between E.K. and the controls in left and right areas 44, in right 45, and in interhemispheric asymmetries. We conclude, that the exceptional language competence of E.K. may be related to distinct cytoarchitectonic features in Broca’s region.
\end{lecture}

\origpage{136}{122}
\subsection{Les fondateurs d’une nouvelle science}
À la fin du 18\textsuperscript{e} siècle, les travaux sur le sanskrit de William Jones ont renforcé de manière décisive la thèse d’une parenté entre les langues européennes et le sanskrit. L’étude des langues a alors pris une nouvelle direction : elle ne s’est désormais plus occupée de grammaire générale mais de grammaire \textit{comparée}. À cette grammaire comparée sont associés de grands noms, comme Friedrich Schlegel, Rasmus Rask, Franz Bopp et August Schleicher.

\minititle{\textcircled{\scriptsize 1} Friedrich Schlegel}
Né en Allemagne, \textbf{Friedrich Schlegel} (1772-1829) a appris le persan et le sanskrit auprès de l’orientaliste français \textbf{Antoine-Léonard Chézy} (1773-1832). Cela a permis à Schlegel d’établir la parenté entre le grec, le latin, le sanskrit, le perse et l’allemand.

\minititle{1. \textit{Sur la langue et la sagesse des Indiens} (1808)}
Le premier livre à développer une idée de typologie linguistique a été publié par Frédéric Schlegel en 1808 : il s’agit de \textit{Über die Sprache und Weisheit der Indier} (« sur la langue et la sagesse des Indiens »), qui a été traduit en français en 1837. L’essentiel de ce livre est consacré à l’examen de la langue sanskrite, que Schlegel compare de façon détaillée à d’autres langues de l’Europe. Ce qui lui paraît décisif dans la comparaison des langues est de cesser de se soucier seulement des racines, comme le faisaient les étymologistes, mais d’intégrer la structure grammaticale. Ce terme de « structure » a une longue histoire en linguistique, mais chez Schlegel il a perdu sa connotation architecturale : elle est médicale et fait référence à l’anatomie.

\minititle{2. Comparatisme et anatomie}
Frédéric Schlegel souligne que la science dont il faut s’inspirer pour faire de la grammaire comparée est l’anatomie comparée :
\begin{lecture}{}
Mais le point décisif qui éclaircira tout, c’est la structure [Struktur] intérieure des langues ou la grammaire comparée, laquelle nous donnera des solutions toutes nouvelles sur la généalogie des langues, de la même manière que l’anatomie comparée a répandu un grand jour sur l’histoire naturelle plus élevée.
\end{lecture}
Il faut savoir qu’à l’époque de la rédaction de son livre, le naturaliste \textbf{Georges Cuvier} (1769-1832), était en pleine gloire, et venait d’achever la publication de son premier grand-œuvre \textit{Anatomie comparée}. Le parallèle établi par Schlegel entre l’histoire des langues et l’histoire naturelle influencera toute la typologie linguistique, jusqu’à nos jours.

\minititle{\textcircled{\scriptsize 2} Rasmus Rask}
Chercheur original, linguiste et philologue danois \textbf{Rasmus Rask} (1787-1832) a visité l’Islande, la Russie, la Perse et l’Inde afin d’étudier les langues de ces pays. Auteur en 1811 d’une \textit{Grammaire de l’islandais} et de \textit{Recherches sur l’origine de l’ancienne langue nordique ou islandaise} (1818),
\origcont{137}{123}Rask a été le premier à mettre en évidence la parenté existant entre le vieux norrois, le gotique, le lithuanien, le vieux slave, le grec ancien et le latin.

Influencé par le naturalisme de \textbf{Jean-Baptiste de Lamarck}, Rask a fait reculer les limites du comparatisme linguistique en recherchant des explications naturelles à la variation linguistique. Avec lui, la langue est comparée à un organisme vivant qui naît, se développe et se dégrade à travers l’usage qu’en font les hommes. D’où une conception originale du changement : une structure ancienne se désintègre, ouvrant une période parfois longue d’instabilité et de « fermentation », à l’issue de laquelle se dégage une nouvelle structure, plus simple. C’est ainsi que le latin a disparu pour être remplacé par les systèmes romans. On doit également à Rask la notion d’« évolution en spirale » des langues. Le langage humain a des débuts extrêmement simples, puis les langues se complexifient, avant que les exigences du progrès matériel et de l’accroissement des communications humaines viennent ensuite les simplifier. Les langues modernes se rapprochent alors de leur simplicité primitive.

Les travaux de Rask sont loin d’être négligeables mais parce qu’il écrivait en danois, ils ont mis beaucoup de temps à se diffuser en Europe. C’est son contemporain Franz Bopp, qui, à la même époque, va faire émerger le comparatisme comme nouvelle science.

\minititle{\textcircled{\scriptsize 3} Franz Bopp}
Après des études classiques en Allemagne, \textbf{Franz Bopp} (1791-1867) s’est installé à Paris en 1812 pour se consacrer à l’étude du persan, de l’arabe, de l’hébreu et du sanskrit, études.\ednote{原文句末多出孤立的études，语法关系不成立，疑为编辑残留；可删该词及其前逗号。正文仍保留。} En 1816, il publie un travail considéré comme l’acte fondateur de la grammaire comparée : \textit{Le système de conjugaison du sanskrit comparé avec celui des langues grecque, latine, persane et germanique}. Bopp y effectue la première comparaison systématique des langues européennes avec le sanskrit. En 1833, Bopp a publié une volumineuse grammaire comparée dans laquelle il ajouté\ednote{原文如此。此处复合过去时缺助动词，应为il a ajouté。} le lituanien, l’avestique et l’allemand : \textit{La Grammaire comparée des langues indo-européennes, comprenant le sanscrit, le zend, l’arménien, le grec, le latin, le lithuanien, l’ancien slave, le gothique et l’allemand}.

En rapprochant une dizaine de langues et en démontrant leur parenté par la confrontation de leurs déclinaisons conjugaisons,\ednote{原文如此。两名词之间漏连接词，按上下文应为déclinaisons et conjugaisons。} Bopp a posé à la fois la méthode et les principes du comparatisme et de la linguistique historique, rompant ainsi avec la tradition littéraire et philologique évoquée dans l’introduction. Son but étant de mettre en évidence les conditions de fonctionnement de la langue, considérée comme un organisme en évolution.

L’organicisme de Bopp a beaucoup été commenté : selon lui, les langues doivent être considérées comme des corps naturels organiques qui se forment selon des lois définies selon un principe de vie interne, se développent, puis peu à peu dépérissent. Selon ses propres mots, le langage est un « organisme actif, en quelque sorte autonome par rapport à l’homme ». Ce choix de la métaphore biologique est lié, on l’a vu, au succès des idées de la paléontologie, du naturalisme et de l’évolutionnisme, mais il y a une autre raison : le terme d’« organisme » avait à l’époque le sens de « système », idée qui sera reprise avec succès par la linguistique moderne.

\origpage{138}{124}
\begin{lecture}{Lecture : \textit{Grammaire comparée} (extrait)}
Je me propose de donner dans cet ouvrage une description de l’organisme des différentes langues qui sont nommées sur le titre, de comparer entre eux les faits de même nature, d’étudier les lois physiques et mécaniques qui régissent ces idiomes, et de rechercher l’origine des formes qui expriment les rapports grammaticaux. [...] Les origines des formes grammaticales se révèle,\ednote{底本引文如此。主语Les origines为复数，动词应作se révèlent；这里只订语法，不声称已核定所引法译版本。} la plupart du temps, d’elles-mêmes, aussitôt qu’on étend le cercle de ses recherches et qu’on rapproche les unes des autres les langues issues de la même famille, qui, malgré une séparation datant de plusieurs milliers d’années, portent encore la marque irrécusable de leur descendance commune. Les rapports de la langue ancienne de l’Inde avec ses sœurs de l’Europe sont en partie si évidents qu’on ne peut manquer de les apercevoir à première vue ; mais, d’autre part, il y en a de si secrets, de si profondément engagés dans l’organisme grammatical que, pour les découvrir [...] il faut employer toute la rigueur d’une méthode scientifique pour reconnaître et montrer que tant de grammaires diverses n’en formaient qu’une seule dans le principe.
\end{lecture}

\minititle{\textcircled{\scriptsize 4} August Schleicher}
\textbf{August Schleicher} (1821-1868) a été le premier linguiste à s’intéresser, par une sorte de projection dans le passé, à la reconstitution des parlers indo-européens. Il a également été le premier à écrire un texte (une fable) en proto-indo-européen, langue hypothétique (car non attestée directement par des témoignages écrits) et reconstituée. Intitulée « Avis akvāsas ka » (« le mouton et les chevaux »), cette fable a été publiée en 1868 sous le titre allemand \textit{Eine Fabel in indogermanischer Ursprache}, c’est-à-dire « une fable en langue originelle indo germanique ». La voici, accompagnée de sa traduction en français :

\begin{lecture}{}
\textbf{Avis akvāsas ka}\par
Avis, jasmin varnâ na ā ast, dadarka akvams, tam, vāgham garum vaghantam, tam, bhāram magham, tam, manum āku bharantam. Avis akvabhjams ā vavakat : kard aghnutai mai vidanti manum akvams agantam. Akvāsas ā vavakant : krudhi avai, kard aghnutai vividvant-svas : manus patis varnām avisāms karnauti svabhjam gharmam vastram avibhjams ka varnā na asti. Tat kukruvants avis agram ā bhugat.

\textbf{Le Mouton et les Chevaux}\par
Un mouton qui n’avait pas de laine vit des chevaux, l’un tirant une lourde charrette, l’un soutenant une grosse charge, l’un conduisant un homme à toute allure. Le mouton dit aux chevaux : « Cela me fait mal au cœur de voir un homme conduire des chevaux. » Les chevaux dirent : « Écoute, le mouton, cela nous fait mal au cœur de voir ceci : l’homme, le maître, se fabrique un chaud vêtement avec la laine du mouton. Et le mouton n’a pas de laine. » Ayant entendu ceci, le mouton s’enfuit dans la plaine.
\end{lecture}

\origpage{139}{125}
Schleicher n’a pas échappé lui non plus à l’influence des théories naturalistes et évolutionnistes de l’époque. En 1861, il s’est lié d’amitié avec \textbf{Ernst Haeckel} (1834-1919), biologiste et professeur d’anatomie comparée qui avait introduit les théories de Charles Darwin en Allemagne. Ensemble, ils ont beaucoup discuté sur l’évolutionnisme, les sciences naturelles, l’étude historique des langues. Suite à ces échanges, Schleicher a publié en 1863 \textit{La théorie de Darwin et la science du langage}, ouvrage dans lequel il revendique son biologisme linguistique et réaffirme son ancrage dans les sciences naturelles :
\begin{lecture}{}
Les langues sont des organismes naturels qui, en dehors de la volonté humaine et suivant des lois déterminées, naissent, croissent, se développent, vieillissent et meurent : elles manifestent donc, elles aussi, cette série de phénomènes qu’on comprend habituellement sous le nom de vie. La glottique, ou science du langage, est par suite une science naturelle.
\end{lecture}

\section{LE CLASSEMENT GÉNÉTIQUE DES LANGUES}
\subsection{Sept mille langues : comment les classer ?}
\textit{Ethnologue, Languages of the World}, souvent abrégé en « Ethnologue », est une publication qui depuis 1951 inventorie toutes les langues du monde. À partir de données fournies par des linguistes experts de chaque région, sa 17\textsuperscript{e} édition publiée en 2013 recense 7105 langues différentes\origfoot{1}{D’autres sources plus récentes seront présentées dans le chapitre 13.}. Face à cette diversité, la linguistique dispose de deux grands principes de classification :
\begin{itemize}
\item La linguistique comparée, qui cherche à définir des \textit{familles} de langues, c’est-à-dire un ensemble de langues regroupées par des liens de parenté génétique. Elle est diachronique (historique) et généalogique.
\item Linguistique typologique, synchronique, qui a pour but le regroupement des langues en \textit{types} de langues en fonction de certaines caractéristiques structurales communes.
\end{itemize}
Ces deux démarches sont différentes mais leur but reste le même : classifier, c’est-à-dire trouver l’unité parmi la diversité des langues.

\minititle{\textcircled{\scriptsize 1} La notion de famille de langues}
En analysant les milliers de langues parlées dans le monde, les linguistes ont pu établir certains liens de parenté plus ou moins étroits entre des langues différentes mais issues d’un même prototype (du grec \textit{protos}, premier, primitif). Selon cette conception généalogique, les langues ont des « parents », des « descendants », des « sœurs », et l’on appelle « famille linguistique » l’ensemble formé de toutes les langues partageant la même origine. La communauté scientifique admet généralement l’existence d’environ 300 familles de langues.

\origpage{140}{126}
\minititle{\textcircled{\scriptsize 2} La notion de branche}
Ces familles comprennent à leur tour des sous-ensembles, appelés « sous-familles » ou « branches ». Ces branches sont constituées de langues individuelles plus ou moins étroitement apparentées entre elles, qui peuvent être des « langues sœurs » ou « langues cousines ». Ainsi, les langues de la branche romane (français, espagnol, italien, espagnol, etc.)\ednote{原文在同一列举中两次写espagnol，疑为重复排印。删除后一项即可消除重复；无证据确定作者原拟列举哪一种语言，故不擅补其他语名。} diffèrent de celles de la branche germanique (anglais, allemand, néerlandais, danois, etc.) et de la branche slave (russe, polonais, tchèque, slovène, etc.). Pourtant, elles appartiennent toutes à la même famille : la famille indo-européenne, ainsi appelée parce que regroupant un grand nombre de langues en usage depuis l’Inde, le Pakistan, l’Iran, l’Iraq, la Syrie et la Russie, jusqu’à l’Ouest de l’Europe du Portugal à Moscou en passant par l’Islande et la Grèce.

L’établissement de liens de parenté entre les langues ne repose pas toujours sur une langue originelle véritable.\ednote{此处把“未有直接文字见证”和“并非真实祖语”混在一起。紧接的说明实际只支持“祖语未直接见于文献、其具体形式依比较重建”，不能由此推出共同祖语不曾真实存在。应将véritable改为directement attestée，或明确区分历史祖语与重建模型。} Dans certains cas, il s’agit d’hypothèses formulées d’après des analyses comparatives et historiques afin de constituer des ensembles de langues. Nous avons vu avec Schleicher que les linguistes ont reconstitué des langues originelles (des protolangues) qui n’ont jamais été attestées et demeurent donc des langues purement hypothétiques. C’est le cas de l’indo-européen, reconstruit par les linguistes mais pour lequel il n’existe aucun document écrit ne pouvant en confirmer l’authenticité.

\minititle{\textcircled{\scriptsize 3} La notion de superfamille}
Les superfamilles (appelées aussi « macro-familles ») regroupent plusieurs familles de langues déjà établies au moyen des méthodes de la linguistique comparative et qui auraient une origine commune. Estimées entre 10 et 20 selon les auteurs, elles ne font pas toutes l’unanimité et suscitent parfois la controverse. \textbf{Joseph Harold Greenberg} (1915-2001) qui a été directeur du département d’anthropologie de l’Université de Stanford a découvert la moitié de ces macro-familles qui regroupent l’ensemble des langues humaines.\ednote{此句的a découvert把前一句仍属有争议的假说改说成了确证发现；在本段自己的证据层级内，应作a proposé。对“半数”及“涵盖全部人类语言”的数量断言，依据不足，仍待核验。} Parmi ses découvertes, citons la superfamille « amérindienne », présentée en 1987 dans son livre \textit{Language in the Americas}, dont l’aire de répartition s’étend du Canada à l’Amérique du Sud, et la superfamille « eurasiatique » (regroupant les langues indo-européennes, ouraliennes, altaïques, turques et langues mongoles) présentée dans \textit{Indo-European and Its Closest Relatives : The Eurasiatic Language Family} (2000 et 2002).

\minititle{\textcircled{\scriptsize 4} La notion d’isolat linguistique}
Pour rester dans la métaphore familiale et généalogique, on pourrait dire qu’un isolat est une langue « née de parents inconnus ». Dans la terminologie linguistique, on dira qu’il s’agit d’une langue dont on ne peut démontrer la relation génétique avec d’autres langues vivantes.\ednote{定义不应只限于langues vivantes（活语言），还须考虑有文献记录的灭绝语言；若与某种灭绝语言的亲缘关系已经确立，也不能仅因该亲属灭绝就称其在谱系分类上为孤立语。见Glottolog对独立语系与孤立语的区分\ecite{24}。} Certaines langues deviennent des isolats lorsque toutes les langues auxquelles elles sont reliées s’éteignent : ce sont en quelque sorte des langues orphelines. Mais d’autres sont des isolats depuis que leur existence est documentée. C’est le cas par exemple du sumérien dont nous avons parlé dans le premier chapitre, mais c’est aussi le cas du basque, seul isolat encore vivant parmi les langues en Europe, et qui ne fait donc pas partie de la famille indo-européenne. L’origine de la langue basque étant antérieure à la diffusion de l’écriture en Europe, on a proposé des relations génétiques entre le
\origcont{141}{127}basque et un nombre invraisemblable de langues mais après deux cents ans de recherches, aucune piste ne s’est imposée, et le mystère demeure.

\minititle{\textcircled{\scriptsize 5} Les langues sinoxéniques}
Le coréen et le japonais, généralement considérés comme des isolats linguistiques,\ednote{至少“日语是孤立语”不准确：日语与琉球诸语言构成日琉语系（Japonic），不能因为尚未确证其与其他语系的关系就把其中的日语单独算作孤立语。见Glottolog的Japonic分类\ecite{25}。} possèdent la caractéristique d’être des langues sinoxéniques, c’est-à-dire des langues dont le vocabulaire possède un grand nombre de termes issus du chinois mais qui n’ont pas de lien génétique avec la famille sinotibétaine. Ce sont, pourrait-on dire, des langues qui ont « un air de famille » avec le chinois. Le terme « sinoxénique » vient de l’anglais \textit{sino-xenic} inventé dans les années 50 par le linguiste américain \textbf{Samuel Martin} (1924-2009), professeur à l’Université de Yale et auteur de nombreux ouvrages sur le coréen et le japonais. Le coréen, le japonais, et le vietnamien, sont trois langues sinoxéniques. Pour des raisons liées à l’histoire, ces trois pays ont importé le système d’écriture chinois (les sinogrammes), ainsi qu’un nouveau vocabulaire, à l’image des mots scientifiques ayant des racines grecques inventés en France ou en Grande-Bretagne. Aujourd’hui, la plupart des mots que l’on trouve dans le dictionnaire de japonais « \textit{Shinsen-kokugo-jiten} » (新明解国語辞典)\ednote{罗马字题名与汉字题名不对应：“新明解”应读Shinmeikai，而Shinsen通常对应“新選”。应作Shinmeikai kokugo jiten。底本未标所据版次与词汇统计口径，“la plupart”的数量判断尚待核，不据题名订正推定其正确。} sont issus directement ou indirectement de la langue chinoise.

\subsection{La famille indo-européenne}
Les langues indo-européennes forment une famille de langues étroitement apparentées dont les éléments lexicologiques, morphologiques et syntaxiques présentent de telles ressemblances qu’elles peuvent se ramener à une langue originelle hypothétique, le proto-indo-européen. Au nombre d’environ un millier, les langues indo-européennes sont actuellement parlées par près de trois milliards de locuteurs sur la planète et sont généralement subdivisées en huit branches : l’albanais, l’arménien, les langues balto-slaves, les langues celtiques, les langues romanes, les langues germaniques, les langues helléniques, et les langues indo-iraniennes.\ednote{这里的“八支”仅列所取分类下仍有活语言的分支；本章后面另列安纳托利亚语支和吐火罗语支。若把本句当作包括灭绝语支在内的完整列表，会与本章自己的后文冲突，宜补明范围。}

\minititle{\textcircled{\scriptsize 1} La branche romane}
On nomme langues romanes l’ensemble les langues\ednote{原文如此，l’ensemble后漏介词，应为l’ensemble des langues。} issues du latin vulgaire, c’est-à-dire la forme de latin vernaculaire (utilisée pour la communication de tous les jours) en \textit{Romania}, espace géographique désignant l’ancien Empire romain d’Occident et l’Empire Romain d’Orient. Les mots « roman » et « romania » remontent tous deux l’adjectif\ednote{原文如此。remonter à的介词缺失，应为remontent tous deux à l’adjectif。} latin \textit{romanus}, qui exprime l’idée que le roman est une langue issue des Romains, par opposition à d’autres langues introduites ultérieurement dans les territoires de l’Empire, comme le francique (dans le nord de la France), langue des Francs appartenant à la branche germanique, et non à la branche romane.

\minititle{1. Diglossie latine}
Si le roman vient du latin \textit{vulgaire}, les Romains vivaient en fait en situation de diglossie, c’est-à-dire qu’il existait deux variétés de latin : une variété haute (le \textit{sermo urbanus}) et une variété basse (le \textit{sermo plebeius}). C’est du \textit{sermo plebeius} qu’est issu le roman, et, à sa suite, toutes les langues romanes.

\origpage{142}{128}
\minititle{1) Du \textit{sermo urbanus} au \textit{sermo plebeius}}
La langue vernaculaire en \textit{Romania} n’était pas le latin classique des textes littéraires appelé \textit{sermo urbanus} (« langue de la ville ») mais une forme distincte très proche, mais plus libre : le \textit{sermo plebeius} (« langue vulgaire, langue du peuple »). Le latin classique ne se limitait pas à la langue écrite : il était couramment parlé parmi les catégories sociales élevées, bien que ces derniers\ednote{原文如此；若回指les catégories sociales，应为ces dernières；若指这些阶层的成员，则须补明确的人称指称。} trouvassent plus raffiné encore de s’exprimer en grec. Lorsque Jules César a été poignardé par Brutus, il ne lui a pas dit « \textit{Tu quoque, fili} » (en latin, « toi aussi, mon fils ») mais « \textit{Kaì sù, téknon} », ce qui est le même sens, mais en grec.\ednote{此句把传闻写成了确定史实。Suétone《Divus Iulius》82先说Caesar未发出言语，随后才把“{\BBFifthTranslit καὶ σὺ τέκνον}”列作另一些人的传述；宜改为“据一种古代传述……”，不应断言他必说了此句\ecite{26}。} Le \textit{sermo plebeius} était la langue des soldats, des commerçants, du petit peuple, bref, de la \textit{plèbe}. Le « \textit{Satyricon} », célèbre roman de l’écrivain romain \textbf{Pétrone}, qui est considéré comme l’un des premiers romans de l’histoire de la littérature, est un témoignage direct de cette diglossie : selon leur catégorie sociale, les personnages du roman (contemporains de l’empereur Néron) s’y expriment dans une langue plus ou moins proche du modèle classique.

\minititle{2) Valérius Probus : l’« Appendix Probi »}
N’ayant jamais accédé au statut de langue littéraire, le latin vulgaire nous est parvenu par des citations, des inscriptions, des registres, de textes courants,\ednote{原文如此，列举中的de应与前项一致，作des textes courants。} mais aussi à travers les critiques des « Vaugelas » de l’époque, défenseurs du latin littéraire, Parmi les textes qui ont consigné les formes jugées décadentes et fautives, il faut retenir l’« \textit{Appendix Probi} », compilation d’erreurs fréquentes relevées par \textbf{Valérius Probus}, grammairien lui aussi contemporain de Néron, comme Pétrone.\ednote{作者归属及年代待核。底本把该材料归于1世纪的Valerius Probus，却未说明所据版本；题名中的Probi本身不足以完成作者鉴定。校订研究全文尚待核读；文献线索见\ecite{27}。} Ce document, qui énumère les erreurs commis\ednote{原文如此，分词与阴性复数erreurs配合，应作commises。} couramment dans la graphie de divers mots du latin classique, offre un intérêt considérable car il met en évidence les corruptions (évolutions) de la prononciation qui ont conduit du latin aux langues romanes. L’« \textit{Appendix Probi} » a la forme d’une liste de 227 paires de mots réparties en deux colonnes : dans celle de gauche figurent les formes correctes du latin classique ; dans celle de droite figurent les formes présentées comme incorrectes. On peut ainsi se faire une idée des différences existant dans l’Antiquité tardive entre les formes écrites et les formes orales correspondantes. Il est important de noter ici que ce sont les formes fautives du latin vulgaire, et non leur équivalent en latin classique, qui sont à l’origine des mots utilisés dans les langues romanes. Pour l’illustrer, on verra ci-dessous quelques exemples tirés de l’« \textit{Appendix Probi} » :
\begin{itemize}
\item \textit{Calida non calda, masculus non masclus, tabula non tabla, oculus non oclus} : ces exemples illustrent le processus d’amuïssement, c’est-à-dire la disparition de voyelles post-toniques brèves. Les mots latins étaient en effet accentués et se prononçaient « \textit{cálida} », « \textit{másculus} », « \textit{tábula} » et « \textit{óculus} ». C’est à cet amuïssement que l’on doit les mots français « chaud », « mâle » (« masle » en ancien français), « table » et « oculaire ».\ednote{末例应区别民间继承与学术借入：oculus经民间音变产生的法语词是œil；oculaire来自拉丁语ocularis，不能作为本段oculus→oclus音变的直接结果\ecite{28}。}
\item \textit{Auctor non autor}. On note ici des réductions de groupes de consonnes. Le groupe /kt/ devient /t/, et \textit{auctor} devient \textit{autor} avant de devenir « auteur » en français. De même, /mt/ est réduit à /t/. C’est le cas du verbe \textit{domitare} devenu \textit{domtare} après amuïssement, puis « dompter ».\ednote{“/mt/变为/t/”与紧接的domtare→dompter例子不符。所示变化是在m与t之间出现p的增音，而不是把m删去；该例不能证明/mt/→/t/。} L’insertion d’un « p » entre le « m » et une consonne occlusive est normale :
\origcont{143}{129}c’est ce qu’on appelle une épenthèse.
\item \textit{Pridem non pride}. Le « m » en fin de mots n’est plus prononcé. Cet amuïssement est, entre autres, à l’origine de la disparition des flexions : les langues romanes, en effet, n’ont plus de déclinaisons, contrairement, par exemple, aux langues germaniques comme l’allemand.\ednote{“罗曼诸语言不再有变格”须限定词类，不能直接推广到整个语言系统。法语je/me等人称代词仍随句法功能改变形式；各罗曼语的名词格情况还须分别核对。}
\end{itemize}

\minititle{2. Naissance de la langue romane}
La première attestation du terme « roman » remonte au célèbre Concile de Tours de mai 813. Lors de ce concile réuni à l’initiative de Charlemagne, une distinction a été faite pour la première fois entre la langue romane et la langue germanique (qualifiée alors de « tudesque »).

\minititle{1) Le concile de Tours (813)}
Lors du concile, les évêques rassemblés par Charlemagne établissent que dans les territoires correspondant à la France et l’Allemagne actuelles, les homélies ne seront plus prononcées en latin mais en « \textit{rusticam Romanam linguam aut Theodiscam, quo facilius cuncti possint intellegere quae dicuntur} », c’est-à-dire dans « langue romane rustique dans la langue tudesque (germanique),\ednote{法译原文如此。拉丁语aut表示“或”，此处漏掉连接关系；应为“en langue romane rustique ou en langue tudesque”。依据为本句并列给出的拉丁原文。} afin que tous puissent plus facilement comprendre ce qui est dit ». Le concile de Tours proclame ainsi la reconnaissance des deux plus grandes composantes linguistiques constituant l’empire de Charlemagne : le monde roman (de tradition latine) et le monde germanique. Ces deux grands ensembles linguistiques ont perduré jusqu’à nos jours, chacun correspondant à une branche linguistique distincte. Il faudra cependant attendre l’an 842 et les \textit{Serments de Strasbourg} pour que le premier texte complet écrit dans en langue romane\ednote{原文dans en重复，应删除dans或改写为dans une langue romane。} soit attesté.

\minititle{2) Les \textit{Serments de Strasbourg} (842)}
Dans l’histoire de France, les \textit{Serments de Strasbourg} du 14 février 842 signent l’alliance militaire de Charles le Chauve et Louis le Germanique, contre leur frère aîné, Lothaire I\textsuperscript{er}, tous trois fils de Louis le Pieux et petit-fils\ednote{原文如此，主语为三人，应为petits-fils。} de Charlemagne. Les \textit{Serments de Strasbourg} précèdent d’un an le Traité de Verdun, d’une importance géopolitique considérable qui verra le découpage de l’ancien empire de Charlemagne. Sur le plan linguistique, ces serments ont également une importance considérable. En effet, Louis le Germanique rédige et prononce son serment en langue romane afin d’être compris des soldats de Charles le Chauve, tandis que Charles le Chauve prononce le sien en langue tudesque (germanique) afin d’être compris par les soldats de Louis. Le texte roman des Serments a une portée philologique et symbolique essentielle puisqu’il constitue pour ainsi dire, « l’acte de naissance de la langue française » dans le cadre d’un accord politique historique. Dans le texte ci-dessous prononcé par Louis le Germanique, la langue romane est à peine séparée du latin vulgaire :
\begin{lecture}{}
Pro Deo amur et pro christian poblo et nostro commun salvament, d’ist di en avant, in quant Deus savir et podir me dunat, si salvarai eo cist meon fradre Karlo et in aiudha et in cadhuna cosa, si cum om per dreit son fradra salvar dift, in o quid il mi altresi fazet, et ab Ludher nul
\origcont{144}{130}plaid nunquam prindrai, qui meon vol cist meon fradre Karle in damno sit.

Pour l’amour de Dieu et pour le peuple chrétien et notre salut commun, à partir d’aujourd’hui, autant que Dieu me donnera savoir et pouvoir, je secourrai ce mien frère Charles par mon aide et en toute chose, comme on doit secourir son frère, selon l’équité, à condition qu’il fasse de même pour moi, et je ne tiendrai jamais avec Lothaire aucun plaid qui, de ma volonté, puisse être dommageable à mon frère Charles.
\end{lecture}
Les troupes de Charles le Chauve promettent, en langue romane :
\begin{lecture}{}
Si Lodhuvigs sagrament, que son fradre Karlo iurat, conservat, et Karlus meos sendra de suo part non lostanit, si io returnar non l’int pois : ne io ne neuls, cui eo returnar int pois, in nulla aiudha contra Lodhuvig nun li iv er.

Si Louis observe le serment qu’il jure à son frère Charles et que Charles, mon seigneur, de son côté, ne le maintient pas, si je ne puis l’en détourner, ni moi ni aucun de ceux que j’en pourrai détourner, nous ne lui serons d’aucune aide contre Louis. (littéralement : ni je ni nul que je puis en détourner, en nulle aide contre Louis ne lui je serai)
\end{lecture}

\minititle{3. Variété des langues romanes}
Si les langues romanes sont toutes issues du latin vulgaire, on peut s’interroger sur leur diversité et leur variété de nos jours. Comment le latin a-t-il donné naissance à l’italien, au français, à l’espagnol, au portugais, au catalan, au roumain, etc.? Une des raisons est que la superficie immense de l’Empire romain et l’absence de norme littéraire et grammaticale ont permis à cette langue vernaculaire de ne pas être figée, ce qui a permis à chaque zone de la \textit{Romania} d’utiliser une « saveur » particulière du latin vulgaire. Par exemple, certaines langues romanes préférèrent utiliser le latin \textit{casa} pour signifier « maison » (« casa » en espagnol, catalan, italien, portugais, roumain), tandis que la langue française a choisi un autre mot latin « \textit{mansio} » de sens identique, qui donnera « \textit{mansion} » en anglais.

\minititle{1) La notion de substrat}
La diversité des langues romanes s’explique aussi la présence\ednote{原文缺介词，应为s’explique aussi par la présence。} dans chaque région de \textit{Romania} d’un substrat local, c’est-à-dire une langue parlée initialement dans une zone, recouverte par une autre (en l’occurrence, le latin vulgaire) qui parvient à laisser quelques traces de sa présence dans la langue d’arrivée. Ainsi, le français moderne a conservé un peu plus d’une centaine de mots de son substrat gaulois. La langue basque, quant à elle, a servi de substrat aux langues ibéro-romanes. Par exemple, « gauche » (qui se dit \textit{sinistrus} en latin classique) se dit « \textit{esquerra} » en catalan, « \textit{izquierda} » en castillan et « \textit{esquerdo} » en portugais, tous trois dérivés du même substrat basque « \textit{ezker} ». Ce terme de substrat a été inventé vers 1875 par le
\origcont{145}{131}linguiste italien \textbf{Graziadio Ascoli} (1829-1907) dans le cadre de ses recherches sur les causes du changement phonétique.

\minititle{2) La notion de superstrat}
En plus des substrats, les superstrats ont aussi joué un rôle prépondérant dans la différenciation des langues romanes. Par superstrat, on entend une langue qui en influence une autre sans toutefois la supplanter. En français, le superstrat francique est très important, et le français actuel compte plusieurs centaines de mots hérités de ce superstrat, notamment le vocabulaire de la guerre et de la vie rurale. Le superstrat arabe a laissé plus de quatre mille mots aux langues castillanes et portugaises.

\minititle{4. Les langues d’oc et langues d’oïl}
En tout début de chapitre, nous avons mentionné le \textit{De l’éloquence vulgaire} de Dante. Dans cet ouvrage, en se basant sur la manière de dire « oui », il distingue trois langues romanes : la langue d’\textit{oïl}, la langue d’\textit{oc} et la langue de \textit{si} (qui correspond à l’italien). Le mot « oïl » vient du gallo-roman \textit{o-il} (« celui-ci »), lequel remonte au latin \textit{hoc et ille}. Le mot « oc », quant à lui, vient du latin \textit{hoc} (« ceci »). Le mot italien \textit{si} vient du latin \textit{sic}, qui signifie « ainsi ».

Langue romane du groupe des langues gallo-romanes, la langue d’oïl s’est développée dans la partie nord de la Gaule, puis dans la partie nord de la France, dans le sud de la Belgique (Belgique romane) et dans les îles Anglo-Normandes. Elle a conservé un substrat celtique important et a subi une grande influence du germanique.

Appelée également « occitan », la langue d’oc est quant à elle parlée dans le tiers sud de la France. Langue régionale la plus parlée en France, son aire linguistique et culturelle est appelée Occitanie (ou Pays d’Oc), elle possède de nombreuses variantes nommées d’après le nom des anciennes provinces : l’auvergnat, le dauphinois, le gascon, le languedocien, le guyennais, le limousin, le provençal. Voilà ce que l’on pouvait dire sur les langues romanes.

\minititle{\textcircled{\scriptsize 2} La branche germanique}
Les langues germaniques constituent l’autre grande branche de la famille des langues indo-européennes. Issues du proto-germanique, elles ont d’abord été parlées par les peuples germaniques qui vivaient initialement aux frontières nord-est de l’Empire romain. Ces langues partagent plusieurs traits uniques, parmi lesquels d’importantes mutations consonantiques. Le philologue allemand \textbf{Jacob Grimm} (un des deux frères Grimm) les a décrites systématiquement en 1822 dans une loi phonétique qui porte son nom : « la loi de Grimm ». Les langues germaniques les plus parlées actuellement sont celles de branche occidentale (anglais, l’allemand, néerlandais) et scandinave (le suédois, le danois et le norvégien), toutes issues de l’indo européen, comme le montrent les deux exemples suivants :
\begin{itemize}
\item L’indoeuropéen \textbf{*phtér} (« père ») a donné naissance à : \textit{fadar} (gotique), \textit{father} (anglais), \textit{fater} (vieux haut allemand), \textit{Vater} (allemand), \textit{vader} (néerlandais), \textit{faðir} (islandais), \textit{fader} (danois et
\origcont{146}{132}suédois), etc.
\item L’indoeuropéen \textbf{*bhére} « porter » a donné naissance à : \textit{baíran} (gotique), \textit{bear} (anglais), \textit{beuren} (néerlandais), \textit{bera} (islandais), \textit{bære} (norvégien), \textit{bära} (suédois), etc.
\end{itemize}
L’astérisque (*) en début de phrase dans les exemples ci-dessus indique qu’il s’agit d’une reconstitution, et c’est à \textbf{August Schleicher} que nous devons cette convention d’écriture.

\minititle{\textcircled{\scriptsize 3} La branche celtique}
Les langues celtiques continentales regroupent les langues parlées par les peuples celtes, depuis l’âge de fer jusqu’à l’Antiquité tardive. Le gaulois (autrefois appelé « gallique ») était l’une de ces langues celtiques. Il a été parlé par les peuples gaulois jusqu’au 5\textsuperscript{e} siècle. Les connaissances liées à cette langue sont lacunaires car les Celtes privilégiaient l’oralité et la mémoire pour la transmission des connaissances. Cependant, on a retrouvé plusieurs inscriptions sur des céramiques, des pièces de monnaie et des objets de la vie quotidienne qui attestent de l’existence de l’écriture chez les Gaulois : d’abord rédigées à l’aide de l’alphabet grec, elles ont ensuite été écrites avec l’alphabet latin après la conquête romaine. La pierre dite « de Martialis », découverte en 1839 dans un sanctuaire dédié au dieu gaulois Ucuetis, est écrite en caractères latins et en langue gauloise. Bien que le gaulois soit éteint depuis le 5\textsuperscript{e} siècle, de nombreux mots subsistent dans certaines langues d’Europe, et surtout les noms des lieux et le vocabulaire rural.

Les langues celtiques insulaires, quant à elles, sont parlées encore de nos jours dans les îles Britanniques et en Bretagne. Ces langues insulaires se divisent elles-mêmes en deux sous-groupes nettement distincts : les langues gaéliques, qui comprennent aujourd’hui l’irlandais, l’écossais et le mannois (parlé sur l’île de Man), et les langues brittoniques, qui comprennent le gallois, le cornique (parlé dans les Cornouailles) et le breton. Il est intéressant de noter que le breton est une langue celte alors que le gallo, la langue d’oïl de la Haute-Bretagne, est une langue romane.

\minititle{\textcircled{\scriptsize 4} La branche indo-iranienne}
L’avestique, vieil-iranien de l’est, était autrefois appelé \textit{zend}, d’après le nom que l’indianiste \textbf{Anquetil-Duperron} avait donné à sa traduction du recueil de textes sacrés de l’Iran : le \textit{Zend-Avesta} (1771). Parent éloigné du vieux-perse (vieil-iranien de l’ouest), l’avestique était aussi génétiquement très proche du sanskrit (vieil-indien du nord-ouest). L’avestique s’écrivait de droite à gauche. Le vieux perse, lui, s’écrivait de gauche à droite au moyen d’une graphie proche du l’écriture\ednote{原文如此，应为de l’écriture。} cunéiforme akkadienne. Si l’avestique et le vieux perse sont aujourd’hui éteints, le sanskrit continue d’être utilisé (très marginalement) à la manière du latin en Occident comme langue liturgique, culturelle et académique par quelques milliers de locuteurs. De nos jours, les deux langues indo-iraniennes les plus parlées aujourd’hui sont l’hindi, descendant du sanskrit (en Inde), et le persan, descendant du vieux-perse (en Iran). Dans la famille iranienne, il faut ajouter dans les langues ouest-iraniennes le groupe kurde (Turquie, Syrie, Irak, Iran, Arménie). L’Afghanistan (avec le pachto) et plusieurs États d’Asie Centrale, pratiquent également largement des langues iraniennes. Enfin, la plupart des langues du Pakistan (comme le pendjabi) et du
\origcont{147}{133}Bangladesh (le bengali) appartiennent à la famille indienne.

\minititle{\textcircled{\scriptsize 5} Les branches éteintes}
La famille indo-européenne possède également deux branches éteintes : la branche anatolienne et la branche tokharienne.

\minititle{1. La branche anatolienne}
L’Anatolie, péninsule située à l’extrémité occidentale de l’Asie, désigne aujourd’hui toute la partie asiatique de la Turquie (Anatolie vient du grec \textit{Anatolē}, qui signifie « Orient »). Pour désigner cette région, le terme d’Asie Mineure est encore très utilisé de nos jours. Les langues anatoliennes (le hittite, le louvite, le palaïte, le carien, le pisidien, le sidétique et le lydien) parlées aux IIe et Ier millénaires avant notre ère, sont toutes éteintes aujourd’hui. Si la Mésopotamie a été le berceau de l’écriture, les civilisations anatoliennes ont joué un rôle déterminant dans l’évolution du langage écrit. Les hittites, par exemple, utilisaient deux écritures : l’écriture cunéiforme et les hiéroglyphes. Les royaumes de l’Asie Mineure antique ont ensuite adopté l’alphabet grec pour noter les premiers textes en langue turque.\ednote{此句混淆了古代安纳托利亚语与突厥语。两者不是同一语系，不能把古代小亚细亚诸王国的文字材料称为最早土耳其语文本；宜改为相应的古代安纳托利亚语言，具体所指仍须原著资料确认。与本章后面将le turc列为非印欧语言的说明亦不相容。} Plus tard, l’Empire ottoman a utilisé l’alphabet arabe pour écrire le turc. Enfin, à la révolution turque (1919-1922), il a été décidé d’utiliser l’alphabet latin, plus moderne et plus accessible, pour une transcription totalement phonétique.\ednote{此处采用拉丁字母的改革应改为1928年，而非1919—1922年。所查二级历史资料将第1353号法律的通过日期列为1928年11月1日；法律原件仍待核\ecite{29}。}

\minititle{2. La branche tokharienne : koutchéen et agnéen}
Les langues tokhariennes (ou \textit{agni-kuči}) étaient des langues parlées et écrites durant le premier millénaire de notre ère dans le bassin du Tarim, dans une région qui correspond aujourd’hui au Xinjiang(Chine). Ces langues tokhariennes ont disparu avec l’arrivée des peuples turcophones, en particulier les Ouïghours, au 9\textsuperscript{e} siècle. Sur les manuscrits ramenés du bassin du Tarim par les expéditions européennes et japonaises au début du 20\textsuperscript{e} siècle figurait une langue inconnue, d’abord appelée la « langue I ». L’orientaliste et turcologue allemand \textbf{Friedrich. W. K. Müller} (1863-1930) lui a donné en 1907 le nom de tokharien. En fait, il n’y a pas une seule langue, mais deux langues : le tokharien A et le tokharien B. L’essentiel des manuscrits en tokharien A proviennent de la région de Karachar, et plus précisément des ruines d’un grand complexe monastique qui se trouve à une trentaine de kilomètres au sud-ouest de cette ville, sur le site de Chortchouq. Dans cette région se trouvait autrefois un royaume appelé « Agni » dans les textes sanskrits. Le tokharien A est donc aussi désigné sous le nom d’agnéen. Quant au tokharien B, il est également nommé « koutchéen », du nom de la région de Koutcha dans laquelle ont été trouvés les manuscrits écrits dans cette langue.

\section{LA LINGUISTIQUE TYPOLOGIQUE}
Derrière la diversité des sept mille langues parlées dans le monde se cachent des schémas communs. L’observation de ces schémas permet d’établir une typologie, et c’est justement la tâche de la linguistique typologique de chercher à classer les langues non pas du point de vue de la parenté génétique, mais d’un point de vue structurel.

\origpage{148}{134}
\subsection{Les grands contributeurs}
Lorsque l’on évoque la linguistique typologique, on ne peut pas ne pas évoquer trois figures incontournables : les allemands Auguste Schlegel, Wilhelm von Humboldt et l’américain Joseph Harold Greenberg.

\minititle{\textcircled{\scriptsize 1} Auguste Schlegel : la tripartition des langues}
Dans \textit{Sur la langue et la sagesse des Indiens} (1808), \textbf{Frédéric Schlegel} proposait déjà d’opposer les langues non pas par des considérations étymologiques (de type langue mère et langue fille), mais par des caractéristiques de structure grammaticale. Dix ans plus tard, en 1818, dans ses Observations sur la langue et la littérature provençales son frère \textbf{Auguste Schlegel} a repris l’idée de son frère d’un classement \textit{structural} des langues, tout en la modifiant légèrement :
\begin{lecture}{}
Les langues qui sont parlées encore aujourd’hui et qui ont été parlées jadis chez les différents peuples de notre globe, se divisent en trois classes : les langues sans aucune structure grammaticale, les langues qui emploient des affixes, et les langues à inflexion.
\end{lecture}
Avec Auguste Schlegel, nous dépassons l’opposition binaire de son frère Frédéric pour arriver à une opposition ternaire (trois classes) basée sur gradation historique. En effet, à propos de la première classe, celles des langues « sans aucune structure grammaticale », il écrit :
\begin{lecture}{}
Les langues de la première classe n’ont qu’une seule espèce de mots, incapables de recevoir aucun développement ni aucune modification. On pourrait dire que tous les mots y sont des racines, mais des racines stériles qui ne produisent ni plantes ni arbres. Il n’y a dans ces langues ni déclinaisons, ni conjugaisons, ni mots dérivés, ni mots composés autrement que par simple juxtaposition, et toute la syntaxe consiste à placer les éléments inflexibles du langage les uns à côté des autres. De telles langues doivent présenter de grands obstacles au développement des facultés intellectuelles.
\end{lecture}

\minititle{\textcircled{\scriptsize 2} Wilhelm von Humboldt}
Linguiste, philosophe, diplomate et ministre de l’éducation prussien, \textbf{Wilhelm von Humboldt} (1767-1835) a été le fondateur de l’université de Berlin qui porte aujourd’hui son nom.

\minititle{1. L’hypothèse humboldtienne}
De ses travaux en tant que linguiste, on retient la typologie des langues. On sait qu’en juin 1820, il a exposé devant l’académie de Berlin les bases d’un gigantesque système de comparaison des langues du monde. De ses travaux, on retient aussi l’hypothèse humboldtienne, selon laquelle les langues véhiculent une vision du monde (\textit{Weltansicht}) propre à chaque communauté humaine, un peu comme si nous regardions le monde à travers des lunettes que les langues nous auraient mises. Dans le discours qu’il a prononcé devant l’académie de Berlin le 20 mars 1826, Humboldt s’est exprimé ainsi :
\origpage{149}{135}
\begin{lecture}{}
La grammaire, plus qu’aucune autre partie du langage, est présente invisiblement dans la façon de penser du locuteur, et chacun apporte dans une langue étrangère ses idées grammaticales, et les dépose, si elles sont plus parfaites et accomplies, dans la langue étrangère.
\end{lecture}
Avec cette hypothèse, Humboldt se fait précurseur de la linguistique moderne et de la célèbre hypothèse de Sapir-Whorf que nous présenterons dans le chapitre 14 consacré à la linguistique américaine. Dans un ouvrage posthume publié en 1936,\ednote{原文如此。此书初刊年应为1836，不是1936；见1836年原版书目与数字化文本\ecite{30}。} Über die Verschiedenheit des menschlichen Sprachbaus und seinen Einfluss auf die geistige Entwicklung des Menschengeschlechts, traduit en français sous le titre « Introduction à l’œuvre sur le kavi », Humboldt formule ainsi cette hypothèse :
\begin{lecture}{}
La pensée, cependant, ne dépend pas seulement de la langue en général, mais, dans une certaine mesure, aussi de chaque langue individuelle déterminée.
\end{lecture}
Humboldt distingue ici le rôle du langage en tant que faculté intellectuelle de la langue particulière pratiquée par le locuteur. Il soutient enfin en outre l’existence de formes grammaticales et de lois universelles régissant l’expression linguistique. Noam Chomsky s’inspirera de cette idée d’universaux linguistiques et retiendra que toute langue exprime dans des structures grammaticales en apparence différentes le même entendement universel.

\minititle{2. Le chinois, « langue sans grammaire » ?}
S’appuyant sur une typologie des langues amorcée par les frères Schlegel, Humboldt est convaincu que les langues indo-européennes, qui disposent d’une grammaire et de flexions, constituent un modèle sensiblement supérieur aux autres langues. Dans cette optique, Humboldt présente le chinois comme « une langue presque entièrement dépourvue de grammaire au sens habituel du mot » et ajoute :
\begin{lecture}{}
Les rapports grammaticaux sont désignés uniquement par la position ou par des mots séparés, et que le lecteur a souvent en charge de deviner à partir du contexte s’il doit prendre un mot pour un substantif, un adjectif, un verbe ou une particule.
\end{lecture}
Dans les deux phrases « 你讨厌！» et « 你讨厌语言学么？», « 讨厌 » est-il un substantif, un verbe ou un adjectif ? Seul le contexte permet de le dire.

Les langues qui ne désignent pas spécifiquement les formes, à commencer par le chinois, sont intéressantes en ce qu’elles mettent en œuvre d’autres procédés d’organisation. Cette différence entre les langues « grammaticalement formées » et celles qui ne le sont pas, débouche pour Humboldt sur deux classes de langues constituant une « différence absolue », reléguant ainsi le chinois, langue sans flexion et donc sans grammaire, parmi les langues primitives, imparfaites, voire incomplètes.

\origpage{150}{136}
\minititle{1) La réponse d’Abel Rémusat}
La réplique du sinologue \textbf{Jean-Pierre Abel Rémusat} (1788-1832), qui venait de publier ses \textit{Éléments de la grammaire chinoise} (1822) ne s’est pas faite attendre. En 1824, dans un article du \textit{Journal asiatique} dont Rémusat était lui-même secrétaire, il a contesté le lien supposé par Humboldt entre les formes grammaticales d’une langue et la plus ou moins grande facilité du développement des idées. Quelques années plus tard, il reviendra sur son opposition aux idées de Humboldt :
\begin{lecture}{}
Cependant le chinois semblait sous quelque rapport faire exception aux principes de l’auteur et on appela son attention sur ce singulier phénomène d’un peuple qui depuis quatre mille ans possède une littérature florissante sans formes grammaticales.
\end{lecture}
Face à Humboldt, il va défendre l’idée que contrairement aux langues indo-européennes le chinois n’exprime que les relations absolument nécessaires pour la compréhension, et qu’il exprime les relations grammaticales par la place des mots dans la phrase.

\minititle{2) Le « génie de la langue chinoise »}
Face aux arguments de Rémusat, Humboldt va nuancer sa position. Abandonnant l’idée d’une éventuelle déficience propre aux langues qui n’ont pas développé de grammaire, il va s’intéresser aux autres procédés d’organisation utilisés par les langues qui ne désignent pas spécifiquement les formes. Humboldt reconnaît que la grammaire n’est pas absente du chinois, mais implicite, « sous-entendue ».

Dans son allocution à l’Académie de Berlin le 20 mars 1826, Humboldt ne considère plus que la langue chinoise est sans grammaire ou à la grammaire imparfaite --- la grammaire chinoise se sert d’une « autre méthode » :
\begin{lecture}{}
En n’adoptant point le système de la distinction des catégories grammaticales des mots, on est dans la nécessité de se servir d’une autre méthode pour faire connaître la liaison grammaticale des idées [...] La langue chinoise emploie tous les mots dans l’état où ils indiquent l’idée qu’ils expriment, abstraction faite de tout rapport grammatical [...] Dans toutes les langues, le sens du contexte doit plus ou moins venir à l’appui de la grammaire. Dans la langue chinoise, le sens du contexte est la base de l’intelligence, et la construction grammaticale doit souvent en être déduite. Le verbe même n’est reconnaissable qu’à son sens verbal. La méthode usitée dans les langues classiques, de faire précéder du travail grammatical et de l’examen de la construction, la recherche des mots dans le dictionnaire, n’est jamais applicable à la langue chinoise. C’est toujours par la signification des mots qu’il faut y commencer. Mais dès que cette signification est bien établie, les phrases chinoises ne prêtent plus à l’amphibologie.
\end{lecture}

\origpage{151}{137}
\minititle{\textcircled{\scriptsize 3} Joseph Harold Greenberg}
En prenant la langue comme objet, les linguistes cherchent à mettre au point des modèles pouvant rendre compte du fonctionnement général des langues. Certains typologues cherchent à décrire et caractériser de façon inductive (c’est-à-dire sur la base d’observations) certaines propriétés communes à toutes les langues : ces propriétés sont appelées « invariants linguistiques » ou « universaux du langage ». Si au 19\textsuperscript{e} siècle Humboldt a fait office de pionnier dans ce domaine, à l’époque moderne ce sont les travaux de \textbf{Joseph Harold Greenberg} qui ont été instrumentaux.

\minititle{1. Les invariants linguistiques}
En avril 1961, au cours de la première conférence interdisciplinaire internationale sur le sujet des universaux linguistiques à New York, l’anthropologue Joseph Harold Greenberg a proposé une liste de 45 universaux morphologiques et syntaxiques basés sur de l’étude de 30 langues.\ednote{原文sur de l’étude中de赘余，应为basés sur l’étude de 30 langues。} Greenberg se base sur trois critères principaux : la présence de prépositions ou de postpositions, l’ordre sujet/verbe/objet dans les phrases (qui sera présenté en détail plus bas), et la position des adjectifs qualificatifs par rapport au nom (avant ou après). À côté des universaux morphologiques et syntaxiques qui concernent le signifiant du signe, on a depuis identifié des universaux phonologiques, grammaticaux, et surtout des universaux sémantiques se rapportant au signifié du signe.

\minititle{2. Typologie syntaxique}
À partir des années 1960, certaines approches ont donné une impulsion nouvelle aux études typologiques en orientant les travaux vers la recherche de correspondances structurales entre les langues. Il s’agit de la typologie syntaxique, branche de la typologie linguistique qui classe les langues selon l’ordre dans lequel les mots apparaissent à l’intérieur des phrases. L’un des points de départ a été l’article de Joseph Greenberg, \textit{Some Universals of Grammar with Particular Reference to the Order of Meaningful Elements}(1963), dans lequel il présente une méthode typologique basée essentiellement sur certains faits de syntaxe. Dans un énoncé, l’ordre des éléments peut se réduire à des formules simples articulées en S (sujet) V (verbe) et O (objet) : on peut ainsi classer les langues selon l’ordre dominant de ces constituants dans la phrase.

\minititle{1) Les langues SOV, SVO et VSO}
Les données statistiques ci-dessous sont issues de l’ouvrage \textit{Basic word order : functional principles}(1986) du linguiste anglais \textbf{Russell Tomlin}. Selon lui, trois types sont dominants dans les langues du monde :
\begin{itemize}
\item Les langues \textbf{SOV} (« le chat la souris mange »). Ce type d’ordre des mots concerne 45\,\% des langues : c’est le plus fréquent. On le trouve par exemple en turc, en coréen, en birman. L’exemple ci-dessous nous vient du japonais, :
\end{itemize}
\noindent\textbf{Ex :} « Taro ga inu o mita », soit littéralement : « Taro le chien a vu ».
\begin{itemize}
\item Les langues \textbf{SVO} (« le chat mange la souris ») représentent 42\,\% des langues du monde
\origcont{152}{138}et comprennent la plupart des langues indo-européennes mais aussi des langues d’Asie du sud-est (la langue khmère du Cambodge, le thaï, le lao du Laos et de Birmanie) ou d’Afrique. L’exemple ci-dessous vient du kiswahili, langue parlée au Kenya :
\end{itemize}
\noindent\textbf{Ex :} « Wanafunzi wasoma vitabu », soit littéralement « les élèves lisent les livres ».
\begin{itemize}
\item Les langues \textbf{VSO} (« mange le chat la souris ») représentent 9\,\% des langues du monde. Elles comprennent notamment l’arabe classique et les langues celtiques insulaires. L’exemple ci-dessous nous vient de l’hébreu biblique et correspond à la première phrase de la Genèse :
\end{itemize}
\noindent\textbf{Ex :} « Bereishit bara Elohim et ha-shamayim v’ha-aretz », soit littéralement : au commencement (\textit{Bereshit}) créa (\textit{bara}) Dieu (\textit{Elohim}) le ciel (\textit{ha’shamayim}) et (\textit{ve}) la terre (\textit{ha’aaretz}).

\minititle{2) Les langues VOS, OVS et OSV}
\begin{itemize}
\item Les langues \textbf{VOS} (« mange la souris le chat ») ne représentent que 3\,\% des langues, et comprennent par exemple le fidjien et le malagasy, une langue parlée à Madagascar et d’où vient l’exemple ci-dessous :
\end{itemize}
\noindent\textbf{Ex :} « Manasa lamba amin’ny savony ny lehilahy », littéralement laver (\textit{Manasa}) habits (\textit{lamba}) avec (\textit{amin}) le savon (\textit{ny savony}) l’homme (\textit{ny lehilahy}), c’est-à-dire « l’homme lave les habits avec le savon ».
\begin{itemize}
\item Les langues \textbf{OVS} (« la souris le chat mange »)\ednote{括号内的顺序是O-S-V，不是标题所说的O-V-S。按本章统一以chat为主语、souris为宾语的例法，OVS应写“la souris mange le chat”。} ne représentent qu’1\,\% des langues. On peut \textit{citer le hixkaryana, langue très rare identifiée par le linguiste} \textbf{Desmond Cyril Derbyshire} (1924-2007), et parlée par à peine plus de 500 personnes au Brésil, le long de la rivière Nhamundá.
\end{itemize}
\noindent\textbf{Ex :} « toto yonoye kamara », soit littéralement l’homme (\textit{toto}) a mangé (\textit{yonoye}) le jaguar (\textit{kamara}), ce qui signifie en fait « le jaguar mangé l’homme ».\ednote{原文法语译文缺助动词a，应作“le jaguar a mangé l’homme”。此处保留源例的语序与转写；作为OVS例，首项为宾语、末项为主语。} Attention aux contresens !
\begin{itemize}
\item Les langues \textbf{OSV}. Le \textit{World Atlas of Language Structures} (WALS), qui répertorie 565 langues, mentionne 4 langues concernées par la syntaxe de type OSV, soit moins de 1\,\% des langues.\ednote{此处565没有说明是哪个版本或子样本，不能当作WALS全库规模。查阅的WALS第81章81A表覆盖1376种语言，其中OSV为4种；还单列无主导语序的语言。它与上文Tomlin的百分比并非同一统计样本，不宜直接拼接\ecite{31}。} On peut citer le \textit{xavánte}, langue parlée par le peuple amérindien des Xavántes au Brésil, ou encore le \textit{warao}, un isolat linguistique parlé par les Amérindiens warao du Venezuela et dans le nord-ouest du Guyana, et qui comporte environ 30000 locuteurs. Voici un exemple en \textit{apurinã}, langue parlée par des peuples du bassin de l’Amazone :
\end{itemize}
\noindent\textbf{Ex :} « anana nota apana », soit littéralement : « ananas - je - aller chercher », c’est-à-dire « je vais chercher un ananas ».

Avant de conclure sur ce point, il faut préciser que les choses ne sont pas si figées que cela, et que de nombreuses langues appartiennent à plusieurs types à la fois. Par exemple, si l’allemand est principalement une langue SVO, il faut faire la distinction entre les propositions principales et les
\origcont{153}{139}propositions subordonnées : l’allemand a une structure SVO dans la phrase principale, mais une structure SOV dans la proposition en subordonnée :

\noindent\textbf{Ex :} (SVO) « Ich liebe dich », soit littéralement « je - aime - te/toi ».

\noindent\hspace*{6mm}(SOV) « Weißt du, dass ich dich liebe ? », soit « (sais-tu que) je - te/toi - aime ».

De même, le français, de structure SVO, devient de structure SOV quand le complément d’objet est un pronom :\ednote{应区分具体构式的线性排列与语言的基本语序分类。法语前置宾语弱代词确实可呈S-O-V，但这不把法语改列为基本SOV语言。WALS第81章明确以主语、宾语均为名词短语的陈述句为比较依据\ecite{31}。}

\noindent\textbf{Ex :} (SVO) Je vois une montagne.

\noindent\hspace*{6mm}(SOV) Je la vois.

\subsection{Typologie morphologique}
Après la typologie syntaxique, il est temps à présent de présenter en détail la typologie morphologique. Si nous savons qu’elle est née des travaux des frères Schlegel et Wilhelm von Humboldt, il est intéressant de noter que pour August Schleicher, les langues ont d’abord été isolantes, constituées d’éléments invariables juxtaposés les uns aux autres. Certaines d’entre elles sont ensuite devenues agglutinantes, constituées d’éléments radicaux auxquels se sont ajoutées des marques grammaticales collées les unes à la suite des autres. Enfin, certaines langues sont devenues flexionnelles, formées d’un radical et de flexions (déclinaisons et conjugaisons). Nous retrouvons ici les idées d’« évolution » et de « perfectionnement » si chères à Humboldt : pour Schleicher, les langues indo-européennes marqueraient ainsi une sorte d’apogée.

\minititle{\textcircled{\scriptsize 1} Les langues isolantes}
Les langues isolantes (appelées aussi analytiques) sont constituées de mots (morphèmes) invariables qui ne connaissent ni conjugaisons, ni déclinaisons. Chaque morphème individuel porte une signification générale : les nuances ou des précisions peuvent être exprimées par l’ajout d’autres mots. Les catégories grammaticales et l’actualisation du sens se font à travers la syntaxe (l’ordre des mots), par l’adjonction de particules, par le contexte par le ton choisi, la plupart d’entre elles étant des langues tonales. Les langues isolantes appartiennent essentiellement à la famille \textit{môn-khmer} orientale (un rameau de la branche \textit{môn-khmer}) et à la famille sino-tibétaine. Le vietnamien, le khmer, le thaï, le laotien, le birman, le chinois, le tibétain, etc. sont des langues isolantes et, à l’exception notable du khmer, tonales. Prenons un exemple tiré du vietnamien :

\noindent\textbf{Ex :} « Khi tôi dèn nhà ban tôi, chung tôi batdau lâm bài »,\ednote{底本越南语例句漏、错多处附加符号。按所给法译拟校为“Khi tôi đến nhà bạn tôi, chúng tôi bắt đầu làm bài”。其中bắt đầu须分词，đến、bạn、chúng、làm的字形不能省略或以其他声调替代；译文所添过去体信息未由本句单独标记。此为语言形式拟校，非原版引文比对。} soit « Quand je suis arrivé chez mon ami nous avons commencé à faire les leçons ».

Le morphème invariable \textit{tôi} apparaît trois fois, à chaque fois sous la même forme inchangée bien que sa fonction grammaticale et la place dans la phrase soient différentes à chaque fois.
\begin{itemize}
\item En début de phrase et précédant le verbe, il indique la première personne du singulier et la fonction sujet, comme dans :
\end{itemize}
\origpage{154}{140}
\noindent\textbf{Ex :} Tôi muốn cà phê (je - veux - un café)
\begin{itemize}
\item Précédé d’un substantif, \textit{tôi} exprime la relation d’appartenance (\textit{ban tôi}, « ami-je », mon ami). En français, en revanche, la relation possesseur possédé est exprimée par le recours à un syntagme prépositionnel (« la maison \textit{de mon ami} »).
\item Enfin, précédé du morphème de la pluralité \textit{chúng} (« plusieurs »), \textit{chúng tôi} signifie « plusieurs je », c’est-à-dire « nous ». En chinois, langue isolante elle-aussi, il y a une petite différence : le morphème de la pluralité « 们 » suit le morphème lexical mais ne précède pas.\ednote{原文末句缺宾语代词，应为mais ne le précède pas。}
\end{itemize}
\noindent\textbf{Ex :} 我 (je) $\rightarrow$ 我们 (nous).

\minititle{\textcircled{\scriptsize 2} Les langues flexionnelles}
En typologie morphologique, une langue flexionnelle est en une langue\ednote{原文est en une中en赘余，应为est une langue。} dans laquelle les mots changent de forme sonore et/ou visuelle selon le rapport grammatical qu’ils entretiennent avec d’autres mots dans une phrase. Dans une langue flexionnelle, contrairement une langue isolante,\ednote{原文缺介词à，应为contrairement à une langue isolante。} de nombreux mots sont variables et changent de forme selon le contexte d’usage. Ce changement de forme est appelé « flexion » et les formes faisant l’objet de la flexion sont dites « fléchies ». Dans certaines langues (latin, allemand, grec, russe, etc.) les noms, les adjectifs et les pronoms sont fléchis. Ces flexions sont appelées « déclinaisons ». Les flexions du verbe, quant à elles, correspondent aux conjugaisons, véritable cauchemar des étudiants en français. On distingue deux catégories de langues flexionnelles :
\begin{itemize}
\item Les langues agglutinantes : dans ces langues, les flexions se font par ajout successif de morphèmes exprimant chacun un trait grammatical (pluralité, temps, possession, mouvement, etc.).
\item Les langues synthétiques (ou fusionnelles) : dans ces langues, les flexions se font aussi par ajout d’une désinence à un radical, mais les différents constituants de la flexion ne sont généralement pas distincts La quasi-totalité des langues indo-européennes sont des langues synthétiques.\ednote{本章把synthétique与fusionnel等同，容易混淆两个分类维度：综合程度涉及一个词包含多少语素，融合性涉及语素边界与形式的可分解程度。黏着结构也可以高度综合，故“黏着／综合”不是严格对立；本段拟比较的是“黏着／融合”\ecite{32}。}
\end{itemize}

\minititle{1. Les langues agglutinantes}
Le terme de langue agglutinante a été créé en 1836 par Wilhelm von Humboldt. Il vient du latin \textit{agglutinare}, qui signifie « coller ensemble ». Selon cette définition, une langue agglutinante est donc une langue dans laquelle les traits grammaticaux sont constitués par l’assemblage de morphèmes dont la forme est quasiment invariable. Le caractère agglutinant consiste dans le fait que ces morphèmes sont collés les uns à la suite des autres, un peu comme les wagons d’un train. Dans les langues agglutinantes, les limites entre chaque morphème restent parfaitement distinctes, chaque « wagon » correspondant à un seul trait grammatical. Parmi les langues agglutinantes figurent notamment les trois isolats que sont le coréen, le japonais, le basque,\ednote{日语的孤立语分类问题见原书127页对应校注及\ecite{25}；形态类型相似不构成谱系孤立或同源的证据。} les langues ouraliennes (l’estonien, le finnois, le hongrois), les langues de la famille dravidiennes du sud de l’Inde (le tamoul), les langues altaïques (le turc, le mongol), ainsi que le kiswahili :

\origpage{155}{141}
\minititle{1) Le \textit{kiswahili}}
Les langues nigéro-congolaises (ou langues Niger–Congo) constituent la plus étendue des familles de langues africaines, tant en répartition géographique qu’en nombre de locuteurs. C’est Joseph Harold Greenberg qui a travaillé à la classification et à la typologie linguistique de ces langues. La très grande majorité des langues d’Afrique subsaharienne appartiennent à cette famille qui comptabilise 1526 langues, soit 21\,\% des langues de la planète. Elles sont parlées par 459 millions de locuteurs, soit 7\,\% de la population mondiale, La famille des langues bantoues forme un sous-ensemble de cette grande famille et regroupe environ 400 langues parlées dans une vingtaine de pays de la moitié sud de l’Afrique. Parmi elles, on trouve les langues \textit{swahili}, dont la variété la plus utilisée et populaire est le \textit{kiswahili}, langue nationale en Tanzanie, au Kenya et en Ouganda, Voici quelques exemples d’agglutination en langue \textit{kiswhili} :\ednote{原文此处漏a，应为kiswahili。本段有关swahili与kiswahili关系的说明另列待核。}

\noindent\textbf{Ex :} \textit{wametulipa} (« ils nous ont payés ») se décompose comme suit :

\noindent\textit{wa} (« ils ») + \textit{me} (morphème du passé) + \textit{tu} (« nous ») + \textit{lipa} (« payer »)

\noindent\textbf{Ex :} \textit{atanipenda} (« il m’aimera ») se décompose comme suit :

\noindent\textit{a} (« il ») + \textit{ta} (morphème du futur) + \textit{ni} (« me ») + \textit{penda} (« aimer »)

Vous trouverez un corpus complet de kiswahili à analyser dans les exercices en fin de chapitre.

\minititle{2) Le finnois}
Le finnois est une langue de la branche fennique des langues ouraliennes. Seuls deux pays au monde ont une langue fennique comme langue officielle : la Finlande (avec le finnois) et l’Estonie (avec l’estonien). Voici un exemple d’agglutination en finnois :

\noindent\textbf{Ex :} taloissani (« Dans mes maisons »), se décompose comme suit :

\noindent\textit{talo} (« maison ») + « \textit{i} » (pluriel) + « \textit{ssa} » (inessif, « dans »), « \textit{ni} » (possessif)

\minititle{3) Le tamoul}
Les langues dravidiennes forment une famille d’une trentaine de langues, essentiellement parlées dans le sud de l’Inde. Parmi les langues dravidiennes, les plus utilisées sont le \textit{télougou} et le \textit{tamoul}, l’une des plus anciennes langues du monde apparue il y a plus de 2000 ans. Parlé dans l’État du Tamil Nadu en Inde du sud, le tamoul a également le statut de langue officielle à Singapour et au Sri Lanka. Langue SOV possédant sept cas, chaque trait grammatical est indiqué par un morphème propre.

Le premier cas, le nominatif, correspond au sujet, est n’est pas fléchi :\ednote{原文est n’est系重复误植，应为et n’est pas fléchi。}

\noindent\textbf{Ex :} « nāy pārkki{\BBFifthTranslit ṟ}atu » (le chien - regarde)

Le deuxième cas, l’accusatif, correspond au COD, indiqué par l’ajout du morphème \textit{ai} :

\origpage{156}{142}
\noindent\textbf{Ex :} « nāN nāy\textbf{ai} pārkki{\BBFifthTranslit ṟ}ēN », littéralement « je - le chien - regarde ».

Le troisième cas, l’instrumental, est indiqué par l’ajout du morphème \textit{ōṭu} :

\noindent\textbf{Ex :} « avaN nāy\textbf{ōṭu} vantāN », littéralement « il - avec le chien - est venu ».

Le quatrième cas, le datif, désigne le COI et la direction avec le morphème \textit{ukku} :

\noindent\textbf{Ex :} « nāy\textbf{ukku} koṭu ! », littéralement « au chien - donne ! ».

Le cinquième cas, l’ablatif, exprime la provenance par l’ajout du morphème \textit{iruntu} :

\noindent\textbf{Ex :} « avaN marattil\textbf{iruntu} vi{\BBFifthTranslit ḻ}untāN », littéralement « il - de l’arbre - est tombé ».

Le sixième cas, le génitif, exprime la possession par l’ajout du morphème \textit{iN} :

\noindent\textbf{Ex :} « nāy\textbf{iN} uṇavu », littéralement « du chien - la nourriture ».

Le septième cas, le complément de lieu, est indiqué par le morphème \textit{il} :

\noindent\textbf{Ex :} « maratt\textbf{il} », littéralement « arbre – dans ».

Quant au pluriel, il est exprimé par l’ajout du morphème \textit{kal}, avec une particularité : il doit s’insérer entre le radical et les autres morphèmes, comme suit :

\noindent\textbf{Ex :} « nāN nāyai pārkki{\BBFifthTranslit ṟ}ēN » (« je regarde le chien »).

\noindent\textbf{Ex :} « nāN nāy\textbf{kal}ai pārkki{\BBFifthTranslit ṟ}ēN » (« je regarde les chiens »).

\minititle{4) Le turc}
Les langues altaïques sont un ensemble de langues parlées en Eurasie, depuis la Turquie et la Moldavie jusqu’à l’Asie de l’Est, en passant par l’Asie centrale, la Sibérie et l’Extrême-Orient russe. Elles rassemblent les langues turques, les langues mongoles et les langues toungouses ainsi que, selon certaines classifications, le coréen et le japonais (considérés comme des isolats par d’autres classifications). Le nom de cette famille de langues dérive de l’Altaï, une chaîne de montagnes d’Asie centrale traversant la Chine, le Kazakhstan, la Mongolie et la Russie. Le turc est une langue SOV particulièrement agglutinante, et ses morphèmes se succèdent selon un ordre établi : pluriel d’abord, puis forme possessive, inessif, etc., comme dans l’exemple ci-dessous :

\noindent\textbf{Ex :} « evlerimde » (« dans mes maisons »), se décompose comme suit :

\noindent\textit{ev} (maison) + \textit{ler} (morphème du pluriel) + \textit{im} (possessif) + \textit{de} (inessif, « dans »).

Comme en tamoul, les noms turcs disposent aussi de déclinaisons, et chaque cas (fonction grammaticale) est indiqué par un morphème spécifique :

\origpage{157}{143}
\begin{center}
\begin{tabularx}{.87\linewidth}{>{\raggedright\arraybackslash}X >{\raggedright\arraybackslash}X >{\raggedright\arraybackslash}X}
\toprule
\textbf{Cas} & \textbf{Turc} & \textbf{Français}\\\midrule
Nominatif & köy & (le) village\\
Accusatif & köy\textbf{ü} & le village\\
Génitif & köy\textbf{ün} & du village\\
Datif & köy\textbf{e} & au village\\
Ablatif & köy\textbf{den} & depuis le village\\
Locatif & köy\textbf{de} & dans le village\\\bottomrule
\end{tabularx}
\end{center}

En turc, un grand nombre d’affixes peuvent venir se greffer les uns après les autres pour créer des « mots » tellement longs que l’on ne sait plus trop s’il s’agit d’un mot ou d’une phrase. Ainsi, le mot le plus long en turc comporte soixante-dix lettres. Il s’agit de : \textit{muvaffakiyetsizleştirici\discretionary{}{}{}leştiriveremeyebileceklerimizdenmiş\discretionary{}{}{}sinizcesine}. Je suis sûr que vous aurez la curiosité d’aller chercher son sens. À côté de cela, le mot le plus long de la langue française, « anticonstitutionnellement », fait bien pâle figure. Mais, à sa décharge, le français n’est pas une langue agglutinante, mais synthétique.

\minititle{2. Les langues synthétiques}
En typologie morphologique, une langue synthétique (appelée aussi « fusionnelle ») est une langue où une seule forme indique plusieurs éléments conceptuels et où les rapports grammaticaux se marquent par des modifications des mots.\ednote{“综合”与“融合”不宜当作同义词；本段对“一种形式表示多个概念成分”的说明属于累积表达。综合程度、语音融合及累积表达须分别界定；参见原书140页校注及\ecite{32}。}

\minititle{1) La notion de paradigme}
On appelle « paradigme » l’ensemble des formes différentes que peut prendre un mot. Le \textit{Dictionnaire de la linguistique} (1974) de Georges Mounin définit ainsi le paradigme :

\begin{lecture}{}
Un paradigme est l’ensemble typique des formes fléchies que prend un morphème lexical combiné avec ses désinences casuelles (pour un nom, un pronom ou un adjectif) ou verbales (pour un verbe), selon le type de rapport qu’il entretient avec les autres constituants de la phrase, selon le nombre, la personne, le temps : on dit déclinaison pour un nom, un pronom ou un adjectif et conjugaison pour un verbe. Ainsi, le paradigme du verbe « être » au présent de l’indicatif est : suis, es, est, sommes, êtes, sont.
\end{lecture}

\minititle{2) Différences entre langues agglutinantes et synthétiques}
Toutes les langues indoeuropéennes sont synthétiques, à l’exception de l’arménien qui est agglutinant.\ednote{此处再次把黏着与综合当成排他类型，不能由“黏着”推出“非综合”。同一种语言还可能因语法范畴不同而同时采用分析和综合手段；将整个印欧语系概括为一种形态类型过强，参见\ecite{32}。} Si langues synthétiques et langues agglutinantes appartiennent toutes deux au groupe des langues flexionnelles, elles se différencient de plusieurs manières.

\begin{itemize}
\item Dans les langues agglutinantes, le radical peut exister seul dans ajout d’affixe.\ednote{原文dans ajout d’affixe应为sans ajout d’affixe。} Ce n’est pas le cas des langues synthétiques. En français, par exemple, le radical « pian- » ne prend son
\origcont{158}{144}
sens qu’après l’ajout des suffixes « -iste(s) » ou « -oter », comme dans « pianiste(s) » et « pianoter ».
\item Dans les langues agglutinantes, les morphèmes sont discernables et les mots composés sont facilement décomposables en éléments de base. Les langues synthétiques présentent au contraire un syncrétisme très poussé : il n’est parfois pas possible de distinguer les limites entre le radical et la désinence. Dans « mes maisons », le morphème « mes » contient deux traits distincts, pluriel et possession, qui ont fusionné et que l’on ne peut plus séparer.
\item Dans les langues agglutinantes, chaque affixe apporte un seul trait grammatical, alors que dans les langues synthétiques, une seule désinence peut apporter plusieurs traits : dans « mes maisons », un seul morphème (« mes ») indique à la fois le nombre (pluriel) et la possession. En turc, il faudrait deux morphèmes (\textit{ler} et \textit{in}) collés l’un à l’autre.\ednote{对应mes“我的（复数领有物）”，此处in应为im；原书142页的ev-ler-im-de已经给出相同分析。in表示第二人称单数领有者，不能在此代替im。}
\item Dans les langues agglutinantes, le morphème est (à de rares exceptions près) toujours invariable. En tamoul, par exemple, il n’existe qu’un seul morphème du pluriel, \textit{kal}, et donc une seule manière d’exprimer le pluriel. Dans les langues synthétiques, les désinences sont variables. Par exemple, en français la désinence du féminin prend plusieurs formes variées : joli $\rightarrow$ jolie ; directeur $\rightarrow$ directrice ; maître $\rightarrow$ maîtresse ; homme $\rightarrow$ hommasse ;\ednote{hommasse不是homme的阴性形式，而是形容“女子外貌或举止粗壮、似男子”的形容词（带贬义）；若举homme的阴性词汇对应，应为femme，且属不同词根的对应，不能分析为加阴性词尾。参见词典释义\ecite{34}。} heureux $\rightarrow$ heureuse ; traître $\rightarrow$ traîtresse, etc. Par ailleurs, en plus d’être variables, les désinences des langues synthétiques comportent aussi de nombreuses irrégularités, véritables cauchemars pour les apprenants. Ainsi, en français, féminin (beau $\rightarrow$ belle ; vieux $\rightarrow$ vieille, etc.) et pluriel (cheval $\rightarrow$ chevaux ; œil $\rightarrow$ yeux ; aïl $\rightarrow$ aulx\ednote{原文aïl误加分音符，应为ail。} etc.) comportent de très nombreuses irrégularités.
\item Enfin, une langue synthétique n’agglutine pas les affixes les uns à la suite des autres. C’est d’ailleurs pour cette raison que les mots comportent un si grand nombre de paradigmes. En effet, puisqu’il n’est pas possible d’aligner sur un axe vertical les traits grammaticaux les uns après les autres, une seule désinence doit contenir en elle-même plusieurs traits grammaticaux. Par exemple, en allemand, dans « Der Mann ist mein Lehrer » (« Le-homme-est-mon-professeur »), un seul morphème grammatical \textit{der} (« le ») contient à lui tout seul quatre traits grammaticaux : défini, singulier, masculin, nominatif (sujet).
\end{itemize}

Pour illustrer de façon plus concrète encore les différences entre langues agglutinantes et langues synthétiques, je vais montrer comment elles font pour exprimer la même idée, et cela à travers deux exemples : l’illatif et l’inessif.

\minititle{3) Illustration : illatif et inessif}
L’illatif (du latin \textit{illere}, « aller ») est le cas grammatical exprimant le mouvement vers un l’intérieur.\ednote{原文vers un l’intérieur多出un，应为vers l’intérieur。illere词源亦可疑，宜参核illatus／inferre；词源仍待核。} L’inessif (du latin \textit{in}, « à l’intérieur », et \textit{essere}, « être ») le cas grammatical indiquant le lieu dans lequel se déroule l’action exprimée par le verbe.\ednote{原文essere不是拉丁语“是”的通常不定式，拉丁语为esse；本句在右括号后还漏est，应为L’inessif (…) est le cas grammatical…。} Les langues agglutinantes et les langues synthétiques ont deux manières différentes de les exprimer.

\origpage{159}{145}
\begin{itemize}
\item Dans les langues agglutinantes, l’inessif est exprimé par un morphème, comme dans les exemples ci-dessous :
\end{itemize}
\begin{center}\small
\begin{tabularx}{\linewidth}{@{}l l l >{\raggedright\arraybackslash}X@{}}
En finnois & : \textit{-ssa} & \textit{talo} (« la maison ») & $\rightarrow$ talo\textbf{ssa} « dans la maison »\\[.5em]
En estonien & : \textit{-s} & \textit{maja} (« la maison ») & $\rightarrow$ maja\textbf{s} « dans la maison »\\[.5em]
En hongrois & : \textit{-ban} & \textit{a ház} (« la maison ») & $\rightarrow$ a ház\textbf{ban} « dans la maison »\\[.5em]
En basque : & : \textit{-an} & \textit{etxe} (« la maison ») & $\rightarrow$ etxe\textbf{an} « dans la maison »\\
\end{tabularx}
\end{center}
\begin{itemize}
\item Dans les langues flexionnelles il s’exprime généralement une préposition\ednote{原文缺介词par，应为il s’exprime généralement par une préposition。此处flexionnelles沿用作者与agglutinantes对比的用法，不能据此认为黏着结构没有屈折。} suivie parfois d’une flexion, comme en allemand (\textit{das} au nominatif $\rightarrow$ \textit{dem} au datif) :
\end{itemize}
\begin{center}\small
\begin{tabularx}{\linewidth}{@{}l l l >{\raggedright\arraybackslash}X@{}}
En espagnol & : \textit{en} & la casa (« la maison ») & $\rightarrow$ \textit{en} la casa\\[.5em]
En allemand & : \textit{in} & \textbf{das} Haus (« la maison ») & $\rightarrow$ \textit{in} \textbf{dem} Haus\\
\end{tabularx}
\end{center}
\begin{itemize}
\item Dans un certain nombre de langues indo-européennes, l’inessif et l’illatif utilisent la même préposition. C’est l’utilisation d’une flexion différente qui permet de distinguer les deux, comme en latin et en allemand par exemple :
\end{itemize}
\begin{adjustwidth}{4mm}{0mm}
En latin :

Illatif : \textit{in} + accusatif\quad « in cubicul\textbf{um} eo »\quad (« je vais dans la chambre »)

Inessif : \textit{in} + datif\ednote{这里应为in + ablatif，不是datif。拉丁语in表示处所时支配夺格，表示进入的方向时支配宾格；cubiculo在此为单数夺格，虽与该词单数与格同形，不能据词尾同形误定句法格。参见Allen与Greenough语法§221\ecite{33}。}\quad « in cubicul\textbf{o} sum »\quad (« je suis dans la chambre »)

En allemand :

Illatif : \textit{in} + accusatif\quad « ich gehe in d\textbf{en} Wagen »\quad (« je vais dans la voiture »)

inessif : \textit{in} + datif\quad « ich bin in d\textbf{em} Wagen »\quad (« je suis dans la voiture »)
\end{adjustwidth}

Enfin, et c’est important, plus une langue est flexionnelle, plus sa syntaxe est souple, puisque les traits grammaticaux et sémantiques ne sont pas exprimés dans l’ordre des mots, mais dans la forme des mots eux-mêmes. Ainsi, en latin, « Paul aime Pierre » peut se dire « \textit{Petrum Paulus amat} », « \textit{Paulus Petrum amat} » ou encore « \textit{Amat Paulus Petrum} » En effet, dans \textit{Paulus}, la flexion \textit{us} indique le nominatif (la fonction sujet). Dans \textit{Petrum}, la flexion \textit{um} indique l’accusatif (le complément d’objet direct). L’ordre des mots relève du style, mais pas de la sémantique.

Nous reviendrons sur ce point dans le chapitre suivant lorsque nous évoquerons un concept cher à Ferdinand de Saussure, celui de rapports syntagmatiques. Le chapitre six, entièrement consacré aux deux géants que sont \textbf{Charles Sanders Peirce et Ferdinand de Saussure}, nous fera entrer de plain-pied dans la linguistique moderne.

\clearpage
\origpage{160}{146}
\section*{Exercices}
\phantomsection\addcontentsline{toc}{section}{Exercices}
\markright{Chapitre 5 — Exercices}
\noindent\textbf{I.\quad Analysez le corpus de \textit{kiswahili} suivant, issu de \textit{Workbook in Descriptive Linguistics} (1955) d’Henry A. Gleason.}

\begin{center}\small
\begin{tabularx}{\linewidth}{@{}>{\raggedright\arraybackslash}X >{\raggedright\arraybackslash}X@{}}
\toprule
1. \textit{atanipenda} : he will like me & 15. \textit{atanipiga} : he will beat me\\[.2em]
2. \textit{atakupenda} : he will like you & 16. \textit{atakupiga} : he will beat you\\[.2em]
3. \textit{atampenda} : he will like him & 17. \textit{atampiga} : he will beat him\\[.2em]
4. \textit{atatupenda} : he will like us & 18. \textit{ananipiga} : he is beating me\\[.2em]
5. \textit{atawapenda} : he will like them & 19. \textit{anakupiga} : he is beating you\\[.2em]
6. \textit{nitakupenda} : I will like you & 20. \textit{anampiga} : he is beating him\\[.2em]
7. \textit{nitampenda} : I will like him & 21. \textit{amenipiga} : he has beaten me\\[.2em]
8. \textit{nitawapenda} : I will like them & 22. \textit{amekupiga} : he has beaten you\\[.2em]
9. \textit{utanipenda} : you will like me & 23. \textit{amempiga} : he has beaten him\\[.2em]
10. \textit{utampenda} : you will like him & 24. \textit{alinipiga} : he beat me\\[.2em]
11. \textit{tutampenda} : we will like him & 25. \textit{alikupiga} : he beat you\\[.2em]
12. \textit{watampenda} : they will like him & 26. \textit{alimpiga} : he beat him\\[.2em]
13. \textit{atakusumbua} : he will annoy you & 27. \textit{wametulipa} : they have paid us\\[.2em]
14. \textit{unamsumbua} : you are annoying him & 28. \textit{tulikulipa} : we paid you\\\bottomrule
\end{tabularx}
\end{center}

\noindent\textbf{II.\quad Quel est le morphème associé aux significations suivantes ?}

\begin{keybox}
\textbf{Sujets :}

Je : \dotfill\hspace*{.55\linewidth}\mbox{}

Tu : \dotfill\hspace*{.55\linewidth}\mbox{}

Il/elle : \dotfill\hspace*{.55\linewidth}\mbox{}

\textbf{Objets :}

Me : \dotfill\hspace*{.55\linewidth}\mbox{}

Te : \dotfill\hspace*{.55\linewidth}\mbox{}

Le/la : \dotfill\hspace*{.55\linewidth}\mbox{}

Nous : \dotfill\hspace*{.55\linewidth}\mbox{}

Vous : \dotfill\hspace*{.55\linewidth}\mbox{}

Eux/elles : \dotfill\hspace*{.55\linewidth}\mbox{}

\textbf{Racines verbales :}

Aimer : \dotfill\hspace*{.55\linewidth}\mbox{}

Ennuyer : \dotfill\hspace*{.55\linewidth}\mbox{}
\end{keybox}
\origpage{161}{147}
\begin{keybox}
Battre : \dotfill\hspace*{.55\linewidth}\mbox{}

Payer : \dotfill\hspace*{.55\linewidth}\mbox{}
\end{keybox}

\noindent\textbf{III. Traduisez en kiswahili :}
\begin{keybox}
Je les ai battus : \dotfill\hspace*{.35\linewidth}\mbox{}

Tu nous as battus : \dotfill\hspace*{.35\linewidth}\mbox{}

Ils m’ont ennuyé : \dotfill\hspace*{.35\linewidth}\mbox{}

Je suis en train de le payer : \dotfill\hspace*{.35\linewidth}\mbox{}
\end{keybox}

\noindent\textbf{IV. Traduisez en français :}
\begin{keybox}
Atanilipa : \dotfill\hspace*{.45\linewidth}\mbox{}

Walikupenda : \dotfill\hspace*{.45\linewidth}\mbox{}

Utawapiga : \dotfill\hspace*{.45\linewidth}\mbox{}

Nimemsumbua : \dotfill\hspace*{.45\linewidth}\mbox{}
\end{keybox}

\section*{Pour aller plus loin}
\phantomsection\addcontentsline{toc}{section}{Pour aller plus loin}
\markright{Chapitre 5 — Pour aller plus loin}
\begin{small}
\setlength{\parindent}{0pt}
\setlength{\parskip}{.55em}
ALIGHIERI (Dante). \textit{De l’éloquence vulgaire}. Paris, Fayard, 2011.

BOPP (Franz). \textit{Grammaire comparée des langues indo-européennes}. Paris, Hachette livre BNF, 1866.

DROIXHE (Daniel). \textit{Souvenirs de Babel}. La reconstruction de l’histoire des langues de la Renaissance aux Lumières, Bruxelles, Académie royale de la langue et de littérature françaises de Belgique, 2007.

GREENBERG (Joseph H.). \textit{Language Universals}. With Special Reference to Feature Hierarchies. De Gruyter Mouton, 2010.

GREENBERG (Joseph H.). \textit{The Languages of Africa}. Indiana University Press, 1966.

HAGÈGE (Claude) Hagège,\ednote{原书书目此处Hagège姓氏重复，应删去右括号后的Hagège。} \textit{La structure des langues}. Coll. « Que sais-je ? », Paris, PUF, 1986.

MEILLET (Antoine). \textit{Introduction à l’étude comparative des langues indo-européennes}. Paris, Hachette livre BNF, 1964.

PHILONENKO (Alexis). \textit{Wilhelm von Humboldt ou L’aurore de la linguistique}. Belles lettres, 2006.

PLARD (Henri). \textit{La correspondance de Marcus Zuerius Boxhornius}. dans La correspondance d’Érasme et l’épistolographie humaniste, Bruxelles, 1985.

\origpage{162}{148}
RAMAT (Paolo). \textit{La typologie linguistique}. Paris, PUF, 1985.

ROUSSEAU (Jean) et THOUARD (Denis). \textit{Lettres édifiantes et curieuses sur la langue chinoise}. Presses universitaires du Septentrion, 1999.

SCHOEBEL (Charles). \textit{La Philologie comparée de l’origine du langage}. Paris, Duprat, 1862.

SCHLEGEL (Friedrich). \textit{Voyage en France}. Montpellier, Grèges, 2003.

TOMLIN (Russell). \textit{Basic word order : functional principles}. Croom Helm, 1986.

TRABANT (Jürgen). \textit{Humboldt ou le sens du langage}. Mardaga, 1995.
\end{small}

\section*{Glossaire}
\phantomsection\addcontentsline{toc}{section}{Glossaire}
\markright{Chapitre 5 — Glossaire}
\gloss{diglossie 双语制，双语现象}{一个社团内同时存在两种语言或语言的两种变体，各用于不同的目的。通常情况下，一种比较标准，称高标准语言（变体），用于政府、新闻媒介、教育和宗教活动；另一种一般为无声望的语言，称低标准语言（变体），用于家里、朋友之间、购物时等。}
\gloss{famille de langues 语系}{据信从一个共同来源发展起来的一组语言。}
\gloss{langue agglutinante 黏着型语言}{语言中的一类，各种词缀加入词根以增加其意义或改变其语法功能。}
\gloss{langue celtique 凯尔特语族}{印欧语系的语族之一。凯尔特诸语言在历史上一度通行于很广的地区。后来由拉丁语逐渐取代了某些凯尔特语言，使它们的通行范围大大缩小，今仅存于不列颠群岛的一些人烟稀少的地区和法国西北部的一个角落。}
\gloss{langue flexionnelle 屈折型语言}{以词形变化说明意义或语法功能变化的一种语言。一般词的基础部分和表示语法功能的部分（如数或时态）之间没有明确区分。}
\gloss{langue germanique 日耳曼语族}{印欧语系的主要语族之一。中欧（德国、奥地利、卢森堡等）、西欧（英国、比利时、荷兰），以及欧洲之外的英语国家和地区。}

\origpage{163}{149}
\gloss{langue indo-européene 印欧语系}{有关联的被认为具有共同始源语（即“原始印欧语”）的若干语言，包括大多数欧洲语言，还包括古印度语梵语和巴利语及印地语一类的语言。\ednote{本词条法文indo-européene少一个n，应为indo-européenne。正文按底本保留。}}
\gloss{langue indo-iranienne 印度 - 伊朗语族}{印欧语系的语族之一。根据语言学家的分析，这个语族的形成与原始印欧人的大迁徙有直接关系。该语族包括印度语支（印度 - 雅利安）与伊朗语支两大语支。}
\gloss{langue isolante 孤立型语言}{词形不变，语法功能由词序和功能词表示的一种语言。}
\gloss{langue romane 罗曼语族}{由拉丁语演化而来的一组语言，包括法语、意大利语、葡萄牙语等。罗曼语言是印欧语系的一个组成部分。}
\gloss{langue synthétique 综合型语言}{有时黏着型语言和屈折型语言也称为综合型语言。\ednote{本条把黏着型也纳入综合型，与原书140、143页正文的排他性分类不一致；综合程度与融合程度应分别说明。保留原释义，并提示读者对照有关校注及\ecite{32}。}}
\gloss{linguistique comparative 比较语言学}{语言学的一个分支，对两种或两种以上的语言作出研究，以比较它们的结构并揭示它们是否有相同之处。}
\gloss{typologie 类型学}{把各种语言分为不同类型的方法。例如，语言可以按其是否为声调语言分类，也可按其最典型的句法结构来分类等。}
